% Content of the blueprint, shared by the web and the pdf version.

\chapter{Introduction}

This blueprint follows \emph{Metastability and phase transition in a social
network model with multiple opinions}, by Felipe Penafiel and Kádmo Laxa
(\href{https://arxiv.org/abs/2607.19651}{arXiv:2607.19651}), and records which
of its statements are formalised in Lean~4 and which are not.

Every numbered statement of the paper appears here.  A green node in the
dependency graph is formalised; a blue one is stated in Lean but its proof is
a \texttt{sorry}.  Chapter~\ref{chap:notes} explains why each unproved
statement is unproved, and records five places where formalising the paper
showed that the written text needs correcting.

Each statement is called here what the paper calls it, and the blueprint
numbers nothing of its own: the node drawn \texttt{prop9} in the dependency
graph is the paper's Proposition~9, \texttt{lem19} its Lemma~19, \texttt{eq1}
its equation~(1), and a reference to any of them prints the paper's own name.
What the paper does \emph{not} state --- a step it leaves implicit, the
plumbing of the sample space, a witness that keeps a vacuous statement from
passing for a theorem --- is numbered A.1, A.2, \dots\ in a sequence of its own
and drawn \texttt{aux-\dots}, so that no name here is ever mistaken for one of
the paper's.

Two files at the root of the repository read this one.  \texttt{STATUS.md} is
generated from it: every statement, its status, and --- for the ones that are
not proved --- which of five kinds of obstruction it has met.
\texttt{FOR-THE-AUTHORS.md} says what each obstruction asks of the authors.

\aux{Remark}
\begin{remark}[What a green node claims, and what it does not]
  \label{note-green}
  A green node means the Lean proof is complete.  On its own it does not say
  that the \emph{paper's} proof has been checked, and the two are not the same
  thing: a statement can be reached by a route the paper does not take, and
  some statements here --- measurability, the plumbing of the sample space ---
  the paper never proves at all.

  The rule this development works under is that a Lean proof \textbf{follows
  the paper's proof} wherever the paper gives one.  Reaching the conclusion
  another way checks nothing about what was written, and checking what was
  written is how \ref{lem19} and~\ref{lem20} came to be flagged.  The exception
  is what the paper leaves implicit --- a step asserted with ``note that'' or
  ``by definition'', and everything measure-theoretic --- which has to be
  supplied, and is then marked as supplied.

  So each proof environment below opens by saying which it is: it follows the
  paper's proof of such and such, it supplies a step the paper asserts, or it
  has no counterpart in the paper.
  Section~\ref{sec:audit} audits every formalised proof on exactly that
  question.
\end{remark}

\section{Conventions}

The network has $N \ge 3$ actors, $\actors = \{1, \dots, N\}$, and $M \ge 2$
opinions, $\opinions = \{1, \dots, M\}$.  A state is an $N \times M$ matrix
$u = (u(a,o))_{a \in \actors, o \in \opinions}$ of social pressures.

\aux{Remark}
\begin{remark}[Scaled coordinates]
  \label{note-scaling}
  The entries of the paper's matrices lie in $\mathbb{Z} + \pressone \mathbb{Z}
  = \pressone \mathbb{Z}$.  Rather than carry a subgroup of $\mathbb{R}$ around,
  the Lean development works throughout with the rescaled matrix
  \[ v(a,o) := (M-1)\, u(a,o) \in \mathbb{Z} . \]
  The rescaling is a bijection, so no statement changes, but the arithmetic
  becomes plain integer arithmetic.  In particular the lattice estimate behind
  \ref{prop8} — that two distinct entries differ by at least $\pressone$ —
  becomes the statement that two distinct integers differ by at least $1$.

  Consequently a Lean statement in this blueprint reads $v$ where the paper
  reads $u$, and the operator of \ref{eq1} adds $M-1$ on the expressed column
  and subtracts $1$ on the others.  The blueprint text below keeps the paper's
  coordinates; the discrepancy is exactly this rescaling.
\end{remark}

\chapter{The model and the main results}

\section{The expression operator and the state space}

\paper{Equation (1)}
\begin{definition}[The expression operator]
  \label{eq1}
  \lean{SocialNetwork.express, SocialNetwork.express_of_ne}
  \leanok
  When actor $a$ expresses opinion $o$, the matrix of social pressures is
  transformed by
  \[
    [\express{a}{o}(u)](b,p) =
    \begin{cases}
      0 & \text{if } b = a, \\
      u(b,p) + 1 & \text{if } b \ne a \text{ and } p = o, \\
      u(b,p) - \pressone & \text{if } b \ne a \text{ and } p \ne o.
    \end{cases}
  \]
\end{definition}

\paper{Equation (2)}
\begin{definition}[The state space]
  \label{eq2}
  \lean{SocialNetwork.IsState, SocialNetwork.IsState.trust_eq_zero}
  \leanok
  \uses{eq1}
  \[
    \Sspace = \Bigl\{ u : \forall a, o,\ u(a,o) \in \mathbb{Z} + \pressone\mathbb{Z};\
      \min_{a} \|u(a,\cdot)\|_\infty = 0;\
      \forall a,\ \textstyle\sum_{o} u(a,o) = 0 \Bigr\}.
  \]
  In scaled coordinates the lattice condition is automatic, so
  \texttt{IsState} carries the two remaining conditions.
\end{definition}

\aux{Lemma}
\begin{lemma}[Conservation of trust]
  \label{aux-trust}
  \lean{SocialNetwork.trust_express, SocialNetwork.trust_express_of_ne, SocialNetwork.trust_express_self}
  \leanok
  \uses{eq1}
  Expressing preserves the trust $T_a(u) = \sum_o u(a,o)$ of every actor except
  the expressing one, whose trust is reset to zero.
\end{lemma}
\begin{proof}
  \leanok
  \textbf{No counterpart in the paper}, which never states trust
  conservation; it is implicit in the row-sum condition of~(2) and in the
  sentence that from any starting position the process remains in $\Sspace$
  once every actor has expressed.
  The increments applied to a row that is not reset cancel: one opinion gains
  $1$ and the remaining $M-1$ opinions lose $\pressone$ each.
\end{proof}

\aux{Lemma}
\begin{lemma}[$\Sspace$ is stable]
  \label{aux-state-stable}
  \lean{SocialNetwork.IsState.express}
  \leanok
  \uses{eq2, aux-trust}
  $\Sspace$ is stable under every $\express{a}{o}$.  This is the claim, made
  just after \ref{eq1}, that from any starting position the process remains in
  $\Sspace$ once every actor has expressed at least once.
\end{lemma}
\begin{proof}
  \leanok
  \textbf{Supplies a step the paper asserts}: ``Indeed, from any starting
  position \dots\ the process will remain in $\Sspace$ forever'', said just
  after~(2) and not argued.  Trust is conserved by \ref{aux-trust}, and the
  expressing actor's row is null, which supplies the second condition.
\end{proof}

\section{Observables}

\aux{Definition}
\begin{definition}[Public opinion and trust, Section~4]
  \label{aux-observables}
  \lean{SocialNetwork.publicOpinion, SocialNetwork.trust}
  \leanok
  \uses{eq2}
  For $o \in \opinions$ the \emph{public opinion} toward $o$ is the total
  pressure exerted by the population, $P_o(u) = \sum_{a \in \actors} u(a,o)$,
  and the \emph{trust} of an actor is $T_a(u) = \sum_{o \in \opinions} u(a,o)$.
\end{definition}

\paper{Remark 3}
\begin{remark}[Trust vanishes on $\Sspace$]
  \label{rem3}
  \lean{SocialNetwork.IsState.trust_eq_zero, SocialNetwork.Bias.Memory.sum_pressure_eq_zero_of_bias_zero}
  \leanok
  \uses{aux-observables, eq2}
  On $\Sspace$ the trust of every actor vanishes identically; this is the
  normalisation built into \ref{eq2}.  The same holds in the biased model
  exactly when $\alpha = 0$, by the identity of \ref{rem1}.
\end{remark}

\section{The jump process}

\paper{Equation (3)}
\begin{definition}[The generator]
  \label{eq3}
  \lean{SocialNetwork.generator, SocialNetwork.jumpRate, SocialNetwork.totalRate, SocialNetwork.jumpWeight}
  \leanok
  \uses{eq1, eq2}
  For $\beta \ge 0$, the polarization coefficient, the process
  $(\Uproc{u}_t)_{t \ge 0}$ is the Markov jump process with generator
  \[
    G f(u) = \sum_{o \in \opinions} \sum_{a \in \actors}
      \exp(\beta u(a,o)) \bigl[ f(\express{a}{o}(u)) - f(u) \bigr],
    \qquad f : \Sspace \to \mathbb{R} \text{ bounded}.
  \]
  We write $\jumprate(u) = \sum_{a,o} \exp(\beta u(a,o))$ for the total jump
  rate, and $(T_n, A_n, O_n)_n$ for the expression times, actors and opinions.
\end{definition}

\paper{Definition 3}
\begin{definition}[The skeleton]
  \label{def3}
  \lean{SocialNetwork.skeletonKernel, SocialNetwork.jumpPMF, SocialNetwork.pathMeasure, SocialNetwork.isMarkovKernel_skeletonKernel, SocialNetwork.skeletonKernel_reachable, SocialNetwork.isState_skeletonKernel, SocialNetwork.skeleton, SocialNetwork.measurable_skeleton, SocialNetwork.returnTime}
  \leanok
  \uses{eq3}
  The skeleton process is $\Uskel{u}_n := \Uproc{u}_{T_n}$.  Its transition
  kernel chooses the pair $(a,o)$ with probability proportional to its rate
  $\exp(\beta u(a,o))$ and moves to $\express{a}{o}(u)$.  For
  $\theta \subset \Sspace$ we write
  $\Rskel{u}(\theta) := \inf\{ n \ge 1 : \Uskel{u}_n \in \theta \}$.
\end{definition}

\aux{Remark}
\begin{remark}[How the process is built in Lean]
  \label{aux-construction}
  \lean{SocialNetwork.ctsPathMeasure, SocialNetwork.jumpTime, SocialNetwork.explosionTime, SocialNetwork.process, SocialNetwork.transitionKernel, SocialNetwork.hittingTimeCts, SocialNetwork.jumpCount_eq_iff, SocialNetwork.measurable_jumpCount, SocialNetwork.measurable_state_ofStepPath, SocialNetwork.measurable_process}
  \leanok
  \uses{def3}
  Mathlib has no theory of continuous-time Markov jump processes on a countable
  state space, but it does not need one: the jump--hold representation is a
  discrete-time chain carrying one extra real coordinate.  Taking the state of
  that chain to be $(A_n, O_n, \text{holding time})$, with the holding time
  exponential of rate $\jumprate$, turns Mathlib's Ionescu--Tulcea theorem
  \texttt{ProbabilityTheory.Kernel.traj} into a construction of the whole
  process.  From it the jump times $T_n$, the process $\Uproc{u}_t$, the
  transition semigroup and the hitting times are plain definitions.

  The chain is driven by the expressed pairs rather than by the matrices, which
  is legitimate because \texttt{Kernel.traj} allows the kernels to depend on the
  whole past.  This buys something essential: a sample point \emph{is} a
  realisation $(A_n, O_n)_n$ and the matrix at time $T_n$ \emph{is} the state
  obtained by replaying it, definitionally.  The deterministic results of
  Chapter~\ref{chap:proofs} therefore transfer to a realisation pointwise, with
  no almost-sure clause anywhere.

  Per \ref{note-green}, the measurability lemmas named above have \textbf{no
  counterpart in the paper}, which does not address measurability;
  \texttt{jumpCount\_eq\_iff} is likewise an artefact of this construction
  rather than a formalisation of anything written there.  The one measurability
  obligation that is not a cylinder argument, \ref{aux-measurable-hitting}, has
  a node of its own.
\end{remark}

\aux{Lemma}
\begin{lemma}[Measurability of the hitting times]
  \label{aux-measurable-hitting}
  \lean{SocialNetwork.measurable_hittingTimeCts,
    SocialNetwork.Bias.measurable_biasedHittingTimeCts,
    SocialNetwork.jumpCount_max_jumpTime,
    SocialNetwork.max_jumpTime_jumpCount_le,
    SocialNetwork.jumpCount_eq_of_not_bddAbove,
    SocialNetwork.hittingCandidates,
    SocialNetwork.sInf_image_eq_hittingCandidates,
    SocialNetwork.measurable_hittingCandidates,
    SocialNetwork.Bias.measurable_stateAfter_ofStepPath}
  \leanok
  \uses{aux-construction}
  $\omega \mapsto \Rhit{u}(\theta)(\omega)$ is measurable, for either process.
\end{lemma}
\begin{proof}
  \leanok
  \textbf{No counterpart in the paper}, which does not address measurability
  anywhere.

  An earlier version of this blueprint said the statement needed an
  almost-sure formulation, or a proof handling the null set on which the path
  misbehaves.  \emph{That was wrong}, and it is worth recording why: the error
  was to take one route to the reduction for the only one.

  The infimum runs over the uncountable family $\{t \ge 0\}$.  The obvious way
  to make it countable is right-continuity of $t \mapsto \Uproc{u}_t(\omega)$,
  and right-continuity does fail --- on realisations whose holding times are
  negative, or whose jump times accumulate from the right, a null set but not
  an empty one.  It is not needed.  Write $S(t) = \{n : T_n \le t\}$, so that
  the jump count at $t$ is $\sup S(t)$, and split on whether $S(t)$ is
  bounded.

  Where it is bounded, $k = \sup S(t)$ both belongs to $S(t)$ and bounds it,
  so $\max(T_k, 0) \le t$ and the jump count there is again $k$: the infimum
  over that level set is \emph{attained}, at a time named by $k$ alone.  Where
  $S(t)$ is unbounded --- the explosion event, on which the supremum returns
  its junk value $0$ --- every larger time has $S$ unbounded too, so the junk
  value persists to the right and every rational above $t$ lies in the same
  level set.

  So the countable family $\{\max(T_k,0)\}_{k \in \mathbb{N}} \cup
  (\mathbb{Q} \cap [0,\infty))$ already meets the infimum, for \emph{every}
  $\omega$, and each of its members is a measurable function of $\omega$.
  Nothing in the argument mentions the holding times, which is why it
  transports to the biased process unchanged: the two hitting times differ only
  in the map sending $k$ to the state after $k$ expressions, and the reduction
  is stated for an arbitrary such map.
\end{proof}

\paper{Theorem 1.1}
\begin{theorem}[Existence]
  \label{thm1-1}
  \lean{SocialNetwork.nonExplosion}
  \leanok
  \uses{eq3, aux-construction}
  For any $\beta \ge 0$ and any $u \in \Sspace$, the jump times satisfy
  $\mathbb{P}(\sup\{T_m : m \ge 1\} = \infty) = 1$.
\end{theorem}
\begin{proof}
  \uses{prop5}
  The jump times are sandwiched between two Poisson processes, splitting the
  expressions into those coming from actors with pressure below $N$ — which
  \ref{prop5} says happen at least once every $N$ steps — and the others.
  \textbf{Not formalised}: Mathlib has no Poisson point process, only the
  Poisson distribution on $\mathbb{N}$.
\end{proof}

\paper{Theorem 1.2}
\begin{theorem}[Uniqueness of the invariant measure]
  \label{thm1-2}
  \lean{SocialNetwork.existsUnique_invariantCts}
  \leanok
  \uses{eq3, aux-construction}
  The process $(\Uproc{u}_t)_t$ has a unique invariant probability measure
  $\muinv$.
\end{theorem}
\begin{proof}
  \uses{thm1-2-skeleton, eq13}
  A uniform Doeblin minorisation gives the skeleton a unique invariant measure
  $\muskel$, which is transferred to continuous time by \ref{eq13}, the
  correspondence being a bijection between the stationary laws of the two
  processes. \textbf{Not formalised}: Mathlib has no Doeblin criterion.
\end{proof}

\paper{Theorem 1.2}
\begin{theorem}[The invariant measure of the skeleton]
  \label{thm1-2-skeleton}
  \lean{SocialNetwork.existsUnique_invariantSkeleton}
  \leanok
  \uses{def3}
  The skeleton process has a unique invariant probability measure $\muskel$
  carried by $\Sspace$.
\end{theorem}
\begin{proof}
  This is the Doeblin step. \textbf{Not formalised}, and it is the keystone:
  \ref{thm1-2}, \ref{thm2}, \ref{thm3}, \ref{thm25}, \ref{thm27} and
  \ref{thm31} all rest on it.
\end{proof}

\paper{Equation (13)}
\begin{lemma}[Transfer to continuous time]
  \label{eq13}
  \lean{SocialNetwork.invariantCts_eq_of_invariantSkeleton}
  \leanok
  \uses{def3, eq3}
  For every $u \in \Sspace$,
  \[
    \muinv(u) = \frac{\muskel(u)/\jumprate(u)}
                     {\sum_{v \in \Sspace} \muskel(v)/\jumprate(v)} .
  \]
\end{lemma}
\begin{proof}
  \textbf{Not formalised}; it needs \ref{thm1-2-skeleton} to be worth stating.
\end{proof}

\section{Ladders, consensus, and the main estimates}

\paper{Definition 1}
\begin{definition}[Ladder sets]
  \label{def1}
  \lean{SocialNetwork.IsLadder, SocialNetwork.ladderSet, SocialNetwork.ladderValues, SocialNetwork.IsLadder.nonneg, SocialNetwork.IsLadder.nonpos}
  \leanok
  \uses{eq2}
  For $o \in \opinions$,
  \[
    \ladder^o = \Bigl\{ u \in \Sspace :
      \{u(1,o), \dots, u(N,o)\} = \{0, \dots, N-1\},\
      u(\cdot,p) = -\tfrac{u(\cdot,o)}{M-1} \ \forall p \ne o \Bigr\},
  \]
  and $\ladder = \bigcup_{o} \ladder^o$.
\end{definition}

\paper{Definition 2}
\begin{definition}[Consensus sets]
  \label{def2}
  \lean{SocialNetwork.IsConsensus, SocialNetwork.consensusSet, SocialNetwork.consensusSetOther}
  \leanok
  \uses{eq2}
  \[
    C^o = \{ u \in \Sspace : u \ne 0,\ u(a,o) \ge 0,\ u(a,p) \le 0
      \ \forall a, \forall p \ne o \},
  \]
  and $C^{-o} = \bigcup_{p \ne o} C^p$.
\end{definition}

\aux{Lemma}
\begin{lemma}[A consensus state has a positive entry]
  \label{aux-consensus-pos}
  \lean{SocialNetwork.IsConsensus.exists_pos}
  \leanok
  \uses{def2}
  For $u \in C^o$ there is an actor with $u(a,o) > 0$.
\end{lemma}
\begin{proof}
  \leanok
  \textbf{Supplies a step the paper asserts}: ``Note that for any $u \in C^o$,
  there exists at least one actor $a$ such that $u(a,o) > 0$'', said just after
  \ref{def2} and not argued.  Were column $o$ identically zero, the vanishing
  row sums would force the whole matrix to vanish, which $C^o$ excludes.
\end{proof}

\aux{Lemma}
\begin{lemma}[$\ladder^o \subset C^o$]
  \label{aux-ladder-consensus}
  \lean{SocialNetwork.IsLadder.isConsensus}
  \leanok
  \uses{def1, def2}
  A ladder supporting $o$ is a consensus state for $o$.
\end{lemma}
\begin{proof}
  \leanok
  \textbf{No counterpart in the paper}: the inclusion is immediate from
  \ref{def1} and~\ref{def2} but is nowhere drawn there.  On a ladder the
  $o$-column is non-negative and the others are the corresponding non-positive
  multiples.
\end{proof}

\paper{Definition 4}
\begin{definition}[Steep ladders]
  \label{def4}
  \lean{SocialNetwork.IsSteepLadder, SocialNetwork.steepLadderSet, SocialNetwork.IsSteepLadder.express}
  \leanok
  \uses{eq2}
  \[
    \steep^o = \bigl\{ u \in \Sspace : u(a,o) \in \mathbb{Z}\ \forall a;\
      0 = u(a_1,o) < \dots < u(a_N,o),\ \{a_1,\dots,a_N\} = \actors;\
      u(a,p) = -\tfrac{u(a,o)}{M-1}\ \forall p \ne o \bigr\} ,
  \]
  and $\steep = \bigcup_o \steep^o$.  These are the states reachable from
  $\ladder^o$ by repeatedly expressing $o$.
\end{definition}

\paper{Remark 5}
\begin{remark}[$\ladder \subset \steep$]
  \label{rem5}
  \lean{SocialNetwork.IsLadder.isSteepLadder, SocialNetwork.ladderSet_subset_steepLadderSet}
  \leanok
  \uses{def1, def4}
  Every ladder is a steep ladder.
\end{remark}
\begin{proof}
  \leanok
  The values $0, 1, \dots, N-1$ are distinct integers with smallest value $0$.
\end{proof}

\paper{Remark 5}
\begin{lemma}[Steep ladders are stable]
  \label{rem5-stable}
  \lean{SocialNetwork.IsSteepLadder.express_of_pos, SocialNetwork.mem_steepLadderSet_of_positivePressure, SocialNetwork.positivePressureEvent, SocialNetwork.measurableSet_positivePressureEvent}
  \leanok
  \uses{def4}
  For $\hat{l} \in \steep$,
  $\{ \Uproc{\hat{l}}_0(A_1,O_1) > 0 \} \subset \{ \Uproc{\hat{l}}_{T_1} \in \steep \}$:
  an expression by a pair carrying strictly positive pressure keeps the state on
  a steep ladder.
\end{lemma}
\begin{proof}
  \leanok
  \textbf{Supplies a step the paper asserts}: Remark~5 writes
  ``Note that $\ladder \subset \steep$ and for any $\hat{l} \in \steep$,
  $\{\dots\} \subset \{\dots\}$'' and gives no argument.
  A positive entry must lie in the supported column, and expressing that opinion
  translates the column off the reset actor, preserving distinctness and the
  minimum $0$.
\end{proof}

\paper{Remark 5}
\begin{lemma}[The bound $\eta$]
  \label{rem5-eta}
  \lean{SocialNetwork.eta, SocialNetwork.eta_le_pathMeasure_positivePressure, SocialNetwork.eta_le_jumpPMF_posFinset, SocialNetwork.posFinset, SocialNetwork.sum_Ico_exp_le_sum_jumpRate, SocialNetwork.sum_jumpRate_compl_le, SocialNetwork.eta_le_div_of_le, SocialNetwork.sum_range_add_le_sum, SocialNetwork.sum_range_eq_one_add_sum_Ico, SocialNetwork.pathMeasure_preimage_zero}
  \leanok
  \uses{def1, def4, rem5-stable}
  For $\hat{l} \in \steep$ and $l \in \ladder$,
  \[
    \mathbb{P}(\Uproc{\hat{l}}_0(A_1,O_1) > 0)
      \ge \mathbb{P}(\Uproc{l}_0(A_1,O_1) > 0) \ge \eta,
    \qquad
    \eta = \frac{\sum_{j=1}^{N-1} e^{\beta j}}
                {\sum_{j=0}^{N-1} e^{\beta j} + MN} .
  \]
\end{lemma}
\begin{proof}
  \leanok
  On a ladder the $N$ pressures for the supported opinion are exactly
  $0, 1, \dots, N-1$, so the numerator collects the rates of the $N-1$ actors
  carrying positive pressure; the remaining entries are non-positive, hence of
  rate at most $1$, and there are at most $MN$ of them.

  \textbf{Supplies a step the paper asserts}: that the worst case over $\steep$
  is attained on $\ladder$.  Remark~5 writes the two inequalities in one line
  and argues neither.  The comparison is
  \texttt{SocialNetwork.sum\_Ico\_exp\_le\_sum\_jumpRate}: on a steep ladder
  the $o$-column is $N$ pairwise distinct non-negative integers one of which is
  $0$, so the positive ones are $N-1$ distinct integers $\ge 1$; being distinct
  they dominate $1, 2, \dots, N-1$ term by term once sorted, and $j \mapsto
  e^{\beta j}$ is increasing.  Every entry outside the $o$-column is
  non-positive by \ref{def4}, so it contributes nothing to the numerator and at
  most $1$ to the denominator.  The domination step is
  \texttt{SocialNetwork.sum\_range\_add\_le\_sum}, the monotone-function form
  of the Mathlib gap recorded below for
  \texttt{SocialNetwork.Bias.sum\_range\_card\_le\_sum}.

\end{proof}

\paper{Remark 5}
\begin{lemma}[The bound $\eta$, iterated]
  \label{rem5-eta-iter}
  \lean{SocialNetwork.eta_pow_le_pathMeasure_steepLadder, SocialNetwork.eta_pow_le_pathMeasure_positiveEvents, SocialNetwork.IsPositiveAt, SocialNetwork.positiveEvents, SocialNetwork.positiveHistory, SocialNetwork.IsSteepLadder.state_of_isPositiveAt, SocialNetwork.eta_le_partialTraj_succ, SocialNetwork.eta_pow_le_historyMeasure}
  \leanok
  \uses{rem5-eta, rem5-stable}
  For $\hat{l} \in \steep$ and any $m$,
  \[
    \mathbb{P}\bigl( \Uskel{\hat{l}}_k \in \steep \ \text{for all } k \le m \bigr)
      \ge \eta^m .
  \]
\end{lemma}
\begin{proof}
  \leanok
  \textbf{Supplies a step the paper asserts}: the iteration itself.  Remark~5
  states the one-step bound; the sketch of Proposition~9 uses the $m$-step one
  --- ``by Remark 5 we show that starting from $l \in \ladder$, a sequence of
  events $\{\Uproc{l}_{T_{j-1}}(A_j,O_j) > 0\}$'' --- without deriving it.

  It is the induction of \ref{prop8} along the finite-horizon kernels, run on
  the event that every expression so far was made from a positive entry.  One
  thing differs, and it is what makes the event move with the state:
  $\zeta_\beta$ bounds the greedy step at \emph{every} matrix, so that
  induction needs nothing of the history it conditions on, whereas $\eta$ holds
  only on $\steep$, so this one has to know that the history has kept the state
  there.  That is exactly what the event says, by \ref{rem5-stable} iterated
  (\texttt{SocialNetwork.IsSteepLadder.state\_of\_isPositiveAt}).
\end{proof}

\paper{Theorem 2}
\begin{theorem}[Concentration of $\muinv$]
  \label{thm2}
  \lean{SocialNetwork.measure_ladderSet_ge, SocialNetwork.tendsto_hittingTime_ladderSet}
  \leanok
  \uses{def1, thm1-2}
  \begin{enumerate}
    \item There is $C > 0$ such that $\muinv(\ladder) \ge 1 - C e^{-\beta/(M-1)}$
      for every $\beta \ge 0$.
    \item For every fixed $\delta > 0$,
      $\sup_{u \in \Sspace \setminus \{0\}}
       \mathbb{P}\bigl(\Rhit{u}(\ladder) > e^{-\frac{\beta}{M-1}(1-\delta)}\bigr) \to 0$
      as $\beta \to +\infty$.
  \end{enumerate}
  The zero matrix must be excluded from the supremum: from $0$ every rate
  equals $1$, so the first expression takes a time of order $1$. \ref{cor11}
  covers it.
\end{theorem}
\begin{proof}
  \uses{prop9, eq13}
  Part 1 follows from \ref{prop9} and the transfer of \ref{eq13}; part 2 from
  \ref{prop7} and \ref{prop8}. \textbf{Not formalised}.
\end{proof}

\paper{Theorem 3}
\begin{theorem}[Metastability]
  \label{thm3}
  \lean{SocialNetwork.metastability}
  \leanok
  \uses{def2}
  There are $\beta_0, C_1 > 0$ and $C_2 \in (0, 1/2)$, depending only on $M$ and
  $N$, such that for $\beta \ge \beta_0$, every $o \in \opinions$ and every
  $u \in C^o$,
  \[
    \sup_{t \ge 0} \left|
      \mathbb{P}\Bigl( \frac{\Rhit{u}(C^{-o})}{\mathbb{E}(\Rhit{u}(C^{-o}))} > t \Bigr)
      - e^{-t} \right| \le C_1 \beta^3 e^{-C_2 \beta},
  \]
  and for any $u, v \in C^o$,
  $\bigl| \mathbb{E}(\Rhit{u}(C^{-o})) / \mathbb{E}(\Rhit{v}(C^{-o})) - 1 \bigr|
   \le C_1 \beta^3 e^{-C_2 \beta}$.
\end{theorem}
\begin{proof}
  \leanok
  \uses{prop12, lem13, lem14, cor15,
    aux-ladder-nonempty}
  Apply \ref{prop12} with the bounds of \ref{lem13} and~\ref{lem14} and
  \ref{cor15}. \emph{Formalised, but neither sorry-free nor axiom-free}: the
  assembly below is complete in Lean, and what it inherits is \texttt{sorryAx}
  from \ref{lem13} and~\ref{lem14} together with the [LM22] axiom of
  \ref{prop12}.

  Section~5.3 checks the four assumptions in half a page; each is one lemma:
  \begin{itemize}
    \item \eqref{eq:h15} at $s_1 = 1$,
      $\varepsilon_1 = 2N^3(M+1)^3 e^{-\beta/(M-1)}$
      (\texttt{measure\_hittingTime\_le\_one}), from part 1 of \ref{lem14}
      through $1 - e^{-x} \le x$;
    \item \eqref{eq:h16} at $s_2 = 2\beta$
      (\texttt{probHittingGT\_ladderOther\_le}), from \ref{lem13};
    \item \eqref{eq:h17} at $\delta = 1/((M+1)N)$
      (\texttt{max\_le\_of\_isCharacteristicTime}), from \ref{cor15};
    \item \eqref{eq:h18} at $\theta = 1/(2(M-1))$
      (\texttt{measure\_hittingTime\_le\_two\_mul}), from part 2 of
      \ref{lem14}.
  \end{itemize}

  Three steps the text leaves implicit had to be supplied.  The paper cites
  \ref{lem13}, which is about $\ladder$, for a condition about $\ladder^o \cup
  C^{-o}$; the two are related by $\ladder \subseteq \ladder^o \cup C^{-o}$,
  since a ladder for $p \ne o$ is a consensus state for $p$.  Steps
  \eqref{eq:h17} and \eqref{eq:h18} both rest on $\sup_{\beta \ge 0} \beta
  e^{-\beta/a} = a e^{-1}$, stated without proof in (20); it follows from $x
  \le e^{x-1}$.  And the threshold $\beta_1$ above which $\varepsilon_1 +
  \varepsilon_2 \le 1/2$ is made explicit, from $e^{-\beta/a} \le a/\beta$.
  Finally, \ref{prop12} fixes $o$ while Theorem~3 quantifies over it, so the
  constants are taken uniform over the finitely many opinions.
\end{proof}

\chapter{Communication bias and the phase transition}

\section{The biased model}

\paper{Equation (5)}
\begin{definition}[The biased operator]
  \label{eq5}
  \lean{SocialNetwork.Bias.Profile.express, SocialNetwork.Bias.Profile.pressure_express, SocialNetwork.Bias.Memory.hear, SocialNetwork.Bias.Memory.reset, SocialNetwork.Bias.Memory.pressure_hear}
  \leanok
  For $\alpha \in \mathbb{R}$, set $\gamma := \pressone - \alpha$ and
  \[
    [\expressa{a}{o}(u)](b,p) =
    \begin{cases}
      0 & \text{if } b = a, \\
      u(b,p) + 1 & \text{if } b \ne a \text{ and } p = o, \\
      u(b,p) - \gamma & \text{if } b \ne a \text{ and } p \ne o.
    \end{cases}
  \]
\end{definition}

\paper{Equation (6)}
\begin{definition}[The biased state space]
  \label{eq6}
  \lean{SocialNetwork.Bias.Profile, SocialNetwork.Bias.IsBiasedState, SocialNetwork.Bias.Memory, SocialNetwork.Bias.Memory.pressure, SocialNetwork.Bias.Memory.heard, SocialNetwork.Bias.IsBiasedState.two_mul_totalHeard_ge, SocialNetwork.Bias.sum_range_card_le_sum, SocialNetwork.Bias.sum_ge_of_injective}
  \leanok
  \uses{eq5}
  An actor that has heard $n_a$ expressions since it last expressed, $c_p$ of
  which concerned opinion $p$, has row $u(a,p) = c_p(1+\gamma) - \gamma n_a$.
  $\Salpha$ is the set of matrices of this form with
  $\min_a \|u(a,\cdot)\|_\infty = 0$ and $\sum_a n_a \ge N(N-1)/2$.

  The Lean development takes the memory profile $(n_a, (c_p)_p)_a$ as the
  primitive datum rather than the matrix, so the row formula holds by
  definition.  See \ref{note-eq6} for a correction to the second condition.

  \texttt{IsBiasedState.two\_mul\_totalHeard\_ge} formalises the paper's own
  argument for the inequality: the $n_a$ are pairwise distinct, whence the
  bound.  The step the paper leaves implicit --- that $N$ distinct naturals sum
  to at least $0 + 1 + \dots + (N-1)$ --- is \texttt{sum\_ge\_of\_injective},
  and has \textbf{no counterpart in the paper} (\ref{note-green}).
\end{definition}

\paper{Remark 1}
\begin{remark}[Arithmetic of a row]
  \label{rem1}
  \lean{SocialNetwork.Bias.IsBiasedState.express, SocialNetwork.Bias.Profile.pressure_mem_addSubgroup, SocialNetwork.Bias.Profile.sum_pressure_eq, SocialNetwork.Bias.IsBiasedState.pressure_ne_zero, SocialNetwork.Bias.Memory.sum_pressure_eq, SocialNetwork.Bias.Memory.pressure_eq_zero_of_heard_eq_zero, SocialNetwork.Bias.Profile.pressure_eq_add_mul, SocialNetwork.Bias.Profile.heard_eq_zero_of_pressure_eq_zero}
  \leanok
  \uses{eq6}
  $\Salpha$ is stable under every $\expressa{a}{o}$.  Moreover
  $\sum_p u(a,p) = n_a (M-1)\alpha$ and
  $u(a,p) = c_p + \gamma(c_p - n_a) \in \mathbb{Z} + \gamma\mathbb{Z}$; a row is
  null if and only if $n_a = 0$ when $(M-1)\alpha \ne 0$, and consequently
  $0 \notin \Salpha$.
\end{remark}
\begin{proof}
  \leanok
  Hearing increases $n_a$ and $c_p$ by one and expressing resets the row to
  $(n_a, c) = (0,0)$, which is a computation on the memory representation.  A
  balanced row with $n_a \ge 1$ has all its entries equal to
  $n_a(M-1)\alpha / M \ne 0$.
\end{proof}

\paper{Remark 2}
\begin{remark}[The degenerate regimes]
  \label{rem2}
  \lean{SocialNetwork.Bias.pressure_lt_pressure_hear_of_neg, SocialNetwork.Bias.le_pressure_hear_of_le_neg_one}
  \leanok
  \uses{eq5}
  For $\alpha > \pressone$, that is $\gamma < 0$, every expression raises the
  pressure of every listener for \emph{every} opinion, so no opinion can acquire
  the negative pressure a consensus requires.  For
  $\alpha \ge 1 + \pressone$, that is $\gamma \le -1$, an expression of $o$
  raises the pressure for every other opinion by at least as much as for $o$
  itself.
\end{remark}
\begin{proof}
  \leanok
  Both are sign computations on the increments of \ref{eq5}.
\end{proof}

\paper{Equation (7)}
\begin{definition}[The biased generator]
  \label{eq7}
  \lean{SocialNetwork.Bias.biasedGenerator, SocialNetwork.Bias.biasedSkeletonKernel, SocialNetwork.Bias.biasedCtsPathMeasure, SocialNetwork.Bias.biasedJumpRate, SocialNetwork.Bias.biasedTotalRate, SocialNetwork.Bias.biasedJumpPMF, SocialNetwork.Bias.biasedPathMeasure, SocialNetwork.Bias.stateAfter, SocialNetwork.Bias.isBiasedState_stateAfter}
  \leanok
  \uses{eq6}
  \[
    \tilde{G} f(u) = \sum_{o \in \opinions} \sum_{a \in \actors}
      \exp(\beta u(a,o)) \bigl[ f(\expressa{a}{o}(u)) - f(u) \bigr].
  \]
\end{definition}

\paper{Equations (8) and (9)}
\begin{definition}[Biased consensus and ladders]
  \label{eq8}
  \lean{SocialNetwork.Bias.IsBiasedConsensus, SocialNetwork.Bias.IsBiasedLadder, SocialNetwork.Bias.IsBiasedSteepLadder, SocialNetwork.Bias.biasedStateSet, SocialNetwork.Bias.biasedLadderSet, SocialNetwork.Bias.biasedConsensusSetOther, SocialNetwork.Bias.biasedSteepLadderSet}
  \leanok
  \uses{eq6}
  For $0 < \alpha < \pressone$,
  \[
    C_\alpha^o = \{ u \in \Salpha : u \ne 0,\ u(a,o) \ge 0,\ u(a,p) \le 0 \},
    \qquad
    \ladder_\alpha^o = \bigl\{ u \in \Salpha :
      \{u(a,o)\}_a = \{0,\dots,N-1\},\ u(\cdot,p) = -\gamma u(\cdot,o) \bigr\},
  \]
  together with $\steep_\alpha$, the biased steep ladders of Appendix~C.
\end{definition}

\aux{Remark}
\begin{remark}[$\ladder_\alpha \subset \steep_\alpha$]
  \label{aux-biased-ladder-steep}
  \lean{SocialNetwork.Bias.IsBiasedLadder.isBiasedSteepLadder, SocialNetwork.Bias.biasedLadderSet_subset_biasedSteepLadderSet, SocialNetwork.Bias.IsBiasedLadder.injective, SocialNetwork.Bias.IsBiasedLadder.exists_zero, SocialNetwork.Bias.IsBiasedLadder.nonneg, SocialNetwork.Bias.IsBiasedLadder.isInt, SocialNetwork.Bias.IsBiasedLadder.exists_eq_natCast}
  \leanok
  \uses{eq8}
  Every biased ladder is a biased steep ladder.
\end{remark}
\begin{proof}
  \leanok
  \textbf{No counterpart in the paper}, which states \ref{rem5} for the
  unbiased model only and transfers it to Appendix~C without comment.  The
  argument is the unbiased one: a column taking each of $0, 1, \dots, N-1$
  exactly once is in particular an injective, non-negative, integer-valued
  column whose minimum is $0$.
\end{proof}

\aux{Lemma}
\begin{lemma}[A biased ladder is a biased consensus state]
  \label{aux-biased-ladder-consensus}
  \lean{SocialNetwork.Bias.IsBiasedLadder.isBiasedConsensus}
  \leanok
  \uses{eq8}
  For $\gamma > 0$ and $N \ge 2$, $\ladder_\alpha^o \subseteq C_\alpha^o$.
\end{lemma}
\begin{proof}
  \leanok
  \textbf{No counterpart in the paper}, exactly as for
  \ref{aux-ladder-consensus}.  The $o$-column is non-negative and takes the
  value $N-1 \ge 1$, so it is not identically zero; the other columns are
  $-\gamma$ times it, hence non-positive.  Appendix~C needs it for the reason
  Section~5.3 needs the unbiased version: \eqref{eq:h16} is about
  $\ladder_\alpha^o \cup C_\alpha^{-o}$ while \ref{lem28} bounds the hitting
  time of $\ladder_\alpha$.
\end{proof}

\paper{Remark 8}
\begin{remark}[The smallest maximum of a non-null row]
  \label{rem8}
  \lean{SocialNetwork.Bias.le_max_pressure, SocialNetwork.Bias.exists_count_ge, SocialNetwork.Bias.pressure_eq_of_bias}
  \leanok
  \uses{eq6}
  For $0 < \alpha < \pressone$ and an actor with $n_a \ge 1$,
  $\max_{p} u(a,p) \ge (M-1)\alpha$, with equality exactly for an actor that has
  heard one expression of each opinion.  Hence the jump rate out of any non-null
  state is at least $e^{\beta (M-1)\alpha}$, which is what replaces
  $e^{\beta/(M-1)}$ in the unbiased model.
\end{remark}
\begin{proof}
  \leanok
  Writing $n_a = (k-1)M + r$ with $k = \lceil n_a/M \rceil$, some opinion has
  been heard at least $k$ times, so
  $\max_p u(a,p) \ge k(1+\gamma) - \gamma n_a = k(M-1)\alpha + \gamma(M-r)
   \ge (M-1)\alpha$.  \textbf{Follows the paper's proof}: the pigeonhole step is
  \texttt{SocialNetwork.Bias.exists\_count\_ge} and the identity is
  \texttt{SocialNetwork.Bias.pressure\_eq\_of\_bias}, written as
  $c_p(1+\gamma) - \gamma n_a = c_p (M-1)\alpha + \gamma(M c_p - n_a)$ so that
  the remainder $r$ need not be introduced: $M c_p - n_a = M - r$ under the
  paper's decomposition.  Both terms are then non-negative, $\gamma > 0$ and
  $M c_p \ge n_a$, and the first is at least $(M-1)\alpha$ because $c_p \ge 1$.
  \textbf{Supplies a step the paper asserts}: that $k \ge 1$, read off $n_a \ge 1$
  without comment; here it follows from $n_a \le M c_p$.
\end{proof}

\section{The phase transition}

\paper{Theorem 4}
\begin{theorem}[Phase transition]
  \label{thm4}
  \lean{SocialNetwork.Bias.biasedAbsorption,
    SocialNetwork.Bias.exists_pos_le_measure_sameFrom,
    SocialNetwork.Bias.sameFrom,
    SocialNetwork.Bias.shiftPath,
    SocialNetwork.Bias.pathMeasure_cylinder,
    SocialNetwork.Bias.pathMeasure_restart,
    SocialNetwork.Bias.measure_inter_shift_eq_sum,
    SocialNetwork.Bias.le_measure_inter_shift,
    SocialNetwork.Bias.measure_inter_shift_le,
    SocialNetwork.Bias.breakHistory,
    SocialNetwork.Bias.biasedProcess, SocialNetwork.Bias.biasedHittingTimeCts}
  \leanok
  \uses{eq7, eq8}
  \begin{enumerate}
    \item For any $\beta > 0$ and any $\alpha < 0$, for any $u \in \Salpha$,
      $\mathbb{P}\bigl( \bigcup_{n \ge 1} \bigcap_{m \ge n}
      \{A_n^{\alpha,\beta,u} = A_m^{\alpha,\beta,u}\} \bigr) = 1$: all but one
      actor eventually stop expressing.
    \item For $0 \le \alpha < \pressone$, adapted versions of Theorems~1, 2
      and~3 hold; these are \ref{thm25}, \ref{thm27} and \ref{thm31}.
  \end{enumerate}
\end{theorem}
\begin{proof}
  \leanok
  \uses{prop17, prop18, thm25,
    thm27, thm31}
  \textbf{Part 1 follows the paper's proof, with its recursion run as a single
  fixed-point step.} \ref{prop17} puts the profile after $N$ expressions in
  $B_N^\alpha$ with probability at least $(NM)^{-N}$, and \ref{prop18} takes
  over from there, so from every $u \in \Salpha$ a single actor performs every
  expression from the $N$-th on with probability at least a constant $c$,
  uniformly in $u$.  The paper then defines the successive failure times
  $\eta_n$, applies the strong Markov property at each and obtains
  $\mathbb{P}(\eta_{n+1} < \infty) \le (1-c)\,\mathbb{P}(\eta_n < \infty)$,
  whence $\lim_n \mathbb{P}(\eta_n < \infty) = 0$.  That limit is exactly $q =
  \sup_{v \in \Salpha} \mathbb{P}_v(\text{no actor is eventually alone})$, and
  the one inequality the induction uses --- the restart at the first failure
  --- gives $q \le (1-c)\,q$ in a single step, which forces $q = 0$ since $c >
  0$.  The estimate, the time it is applied at and the constant are the
  paper's; only the bookkeeping of the recursion is replaced by the fixed point
  it converges to.

  \textbf{Supplies two steps the paper asserts.}  The first is the Markov property itself.
  The paper writes ``the strong Markov property at time $T_N$'', but everything in part~1 is
  about the sequence $(A_n)$, that is about the skeleton, for which $T_N$ is the
  \emph{deterministic} index $N$; so what is used there is the simple Markov property, proved
  here from the exact law of a cylinder.  At the failure times a genuine stopping time does
  appear, and for a discrete-time chain the strong Markov property follows from the simple
  one by decomposing over the countably many values of that time --- which is what the
  failure \emph{events} of the Lean proof do.  The second is that $\{\eta_n < \infty\}$ is
  measurable for the past at $\eta_n$: here it is the statement that each failure event is
  decided by the expressions that precede it.

  \ref{thm16} is not needed for part~1: the statement is about the sequence of
  expressed pairs, and non-explosion is what makes the \emph{continuous-time}
  process well defined.

  Part 2 is \ref{thm25}, \ref{thm27} and~\ref{thm31}, stated separately and
  carrying their own status; Theorem~4 as a whole therefore rests on them.
\end{proof}

\chapter{Proofs}
\label{chap:proofs}

\section{The deterministic layer}

The results of this section are statements about a \emph{realisation} $(A_n,
O_n)_n$, not about its probability.  They are formalised on the type of
realisations and, by \ref{aux-construction}, apply to a sample point of the
process pointwise.

\aux{Definition}
\begin{definition}[Realisations]
  \label{aux-trajectory}
  \lean{SocialNetwork.Trajectory, SocialNetwork.Trajectory.state, SocialNetwork.skeleton, SocialNetwork.Trajectory.state_succ, SocialNetwork.rowSup}
  \leanok
  \uses{eq1}
  A realisation is the sequence $(A_n, O_n)_n$ of expressing actors and
  opinions; $\Uskel{u}_n$ is obtained by iterating $\express{A_n}{O_n}$ from
  $u$.  Note the index convention: in Lean, \texttt{actor 0} is the paper's
  $A_1$.
\end{definition}

\aux{Lemma}
\begin{lemma}[$\Sspace$ along a realisation]
  \label{aux-state-along}
  \lean{SocialNetwork.Trajectory.isState_state, SocialNetwork.isState_skeletonKernel}
  \leanok
  \uses{aux-trajectory, aux-state-stable}
  If $u \in \Sspace$ then $\Uskel{u}_n \in \Sspace$ for every $n$.
\end{lemma}
\begin{proof}
  \leanok
  \textbf{No counterpart in the paper}, which uses this silently throughout
  Section~5.  Induction on $n$ using \ref{aux-state-stable}.
\end{proof}

\paper{Proposition 5}
\begin{proposition}[Low pressure recurs]
  \label{prop5}
  \lean{SocialNetwork.exists_rowSup_actor_lt, SocialNetwork.exists_rowSup_actor_lt_ofPath, SocialNetwork.rowSup_state_le, SocialNetwork.actor_ne_of_rowSup_state_zero, SocialNetwork.actor_ne_actor_of_lt}
  \leanok
  \uses{aux-trajectory, aux-state-along}
  For any $u \in \Sspace$,
  $\inf\{ n \ge 1 : \|\Uskel{u}_{n-1}(A_n,\cdot)\|_\infty < N \} \le N$: among any
  $N$ consecutive expressions, at least one comes from an actor whose social
  pressure on the expressed opinion is below $N$.
\end{proposition}
\begin{proof}
  \leanok
  \textbf{Follows the paper's proof of Proposition~5}, step for step.
  Suppose not.  Since $u \in \Sspace$ some actor $a_0$ has a null row, and a row
  grows by at most $1$ per expression, so $\|\Uskel{u}_m(a_0,\cdot)\|_\infty \le m$;
  hence $A_n \ne a_0$ for $n \le N$.  An actor that has just expressed has a null
  row, so the same argument gives $A_n \ne A_j$ for $j < n \le N$.  That exhibits
  $N+1$ distinct actors in a network of $N$.
\end{proof}

\aux{Definition}
\begin{definition}[The greedy event]
  \label{aux-greedy}
  \lean{SocialNetwork.IsGreedyAt, SocialNetwork.greedyEvent, SocialNetwork.greedyEvents, SocialNetwork.measurableSet_greedyEvent, SocialNetwork.measurableSet_greedyEvents}
  \leanok
  \uses{aux-trajectory}
  For $u \in \Sspace$ and $n \ge 1$,
  \[
    \xi_n^u := \Bigl\{ (A_n, O_n) \in
      \operatorname*{arg\,max}_{(a,o) \in \actors \times \opinions}
      \Uskel{u}_{n-1}(a,o) \Bigr\} :
  \]
  the $n$-th pair is chosen in the most likely way.  It constrains only finitely
  many coordinates, so it is a measurable cylinder.
\end{definition}

\paper{Proposition 6}
\begin{proposition}[Confinement]
  \label{prop6}
  \lean{SocialNetwork.entry_mem_of_greedy, SocialNetwork.entry_mem_of_mem_greedyEvents, SocialNetwork.entry_le_of_greedy, SocialNetwork.neg_le_of_forall_le}
  \leanok
  \uses{aux-greedy, prop5}
  For any $u \in \Sspace$,
  $\bigcap_{j=1}^{N} \xi_j^u \subset \{ -MN < \Uskel{u}_N < N \}$, entrywise.
\end{proposition}
\begin{proof}
  \leanok
  \textbf{Follows the paper's proof of Proposition~6}, in its two cases.  One
  step is supplied: where the paper passes from the bound at the repeat time
  $m$ to ``this implies $\max \Uskel{u}_N < N$'', the Lean proof carries the
  bound over the remaining $N - m$ expressions explicitly.  If the $N$
  expressing actors are distinct, every row was reset within the window and the
  argument of \ref{prop5} bounds the matrix by $N-1$.  Otherwise some actor
  repeats at a step $m$, and greediness at that step caps the whole matrix by
  the small row of the repeating actor.  The lower bound then comes from the
  vanishing row sums.
\end{proof}

\aux{Lemma}
\begin{lemma}[From consensus to a ladder]
  \label{aux-consensus-ladder}
  \lean{SocialNetwork.isLadder_state, SocialNetwork.isLadder_state_of_mem_greedyEvents, SocialNetwork.opinion_eq_of_greedy, SocialNetwork.IsConsensus.express, SocialNetwork.actor_ne_actor_of_greedy, SocialNetwork.state_add_of_hearing, SocialNetwork.opinion_eq_of_mem_greedyEvents}
  \leanok
  \uses{aux-greedy, def2, def1, aux-consensus-pos}
  If $v \in C^o$ and $\bigcap_{j=1}^N \xi_j^v$ occurs then $\Uskel{v}_N \in \ladder^o$.
\end{lemma}
\begin{proof}
  \leanok
  \textbf{No counterpart in the paper.} Proposition~7 closes with ``we consider
  that by definition, if $v \in C^o$ and $\bigcap_{j=1}^N \xi_j^u$ occurs, then
  $\Uskel{u}_N \in \ladder^o$'', and it is not by definition --- see
  \ref{note-prop7}.  The three stages below are the formalisation's own.  In
  $C^o$ a maximising entry lies in column $o$ by \ref{aux-consensus-pos}, and
  expressing $o$ keeps the state in $C^o$, so every one of the $N$ expressions
  expresses $o$.  No actor expresses twice: one that already expressed carries
  too little pressure, while an actor that has not expressed has been
  accumulating pressure for $o$ the whole time.  Each actor therefore expresses
  exactly once, at some step $i$, and afterwards only listens to expressions of
  $o$; its row at time $N$ is exactly $N-1-i$ on column $o$.  As $i$ runs over
  $0, \dots, N-1$ so does $N-1-i$, which is \ref{def1}.
\end{proof}

\paper{Proposition 7}
\begin{proposition}[Reaching a ladder]
  \label{prop7}
  \lean{SocialNetwork.isLadder_state_of_greedy, SocialNetwork.greedyEvents_subset_ladder, SocialNetwork.one_sub_le_pathMeasure_ladder, SocialNetwork.consensusUnion}
  \leanok
  \uses{aux-greedy, def1}
  For any $u \in \Sspace$,
  $\bigcap_{j=1}^{(M+1)N} \xi_j^u \subset \{ \Uskel{u}_{(M+1)N} \in \ladder \}$.
\end{proposition}
\begin{proof}
  \leanok
  \uses{lem19, lem20, aux-consensus-ladder, aux-shift}
  \ref{lem19} needs $\tau(u) \le N+1$ steps to reach $\bigcup_o S^o$,
  \ref{lem20} needs $(M-1)(n(u)+1)-1 \le (M-1)N$ further steps to reach
  $\bigcup_o C^o$, and \ref{aux-consensus-ladder} needs $N$ more, for a total
  of $(M+1)N$.  Composing the three requires restarting the realisation at an
  intermediate time, which is \ref{aux-shift}. \emph{Formalised, but not
  sorry-free}: the assembly is complete in Lean and inherits \texttt{sorryAx}
  from \ref{lem19} and~\ref{lem20} alone.

  \textbf{Supplies a step the paper asserts, and takes a different route
  through one of them.} The written proof carries the process to $\bigcup_o
  S^o$ \emph{at time $N+1$} --- ``so by Lemma~19 we have
  $\bigcap_{j=1}^{N+1}\xi_j^u \subset \{\Uskel{u}_{N+1} \in \bigcup_o S^o\}$''
  --- whereas \ref{lem19} delivers that membership at the first repeat
  $\tau(u)$, which is only $\le N+1$.  Bridging the gap needs $\bigcup_o S^o$
  to be stable under a greedy expression, and the paper neither states nor
  proves that.  It is not obvious either: a greedy expression in $S^o$
  expresses $o$ (\ref{lem20-opening}), but it resets the row of the actor that
  expresses, and that actor may be one of the very witnesses $a_1(u), \dots,
  a_{n(u)}(u)$; new witnesses have to be produced, and the bound on the other
  columns then has to be re-established, which is precisely the bookkeeping of
  \ref{lem20} whose written invariant does not survive (\ref{note-lem20}).

  \ref{lem20} is therefore applied here at $\tau(u)$ itself, where \ref{lem19}
  leaves the process, and no stability of $S^o$ is needed.  The arithmetic is
  the paper's and is unchanged: $\tau(u) + (M-1)(n(u)+1) - 1 \le (N+1) + (M-1)N
  - 1 = MN$.

  The written proof needs a \emph{second} bridge of the same kind, from
  $N+1+(M-1)(n(u)+1)-1$ up to $MN$, and that one is unavoidable --- the last
  stage has to start at $MN$ for the total to be $(M+1)N$.  It is also
  available: it is the paper's own observation, from the last step of this same
  proof, that a greedy expression in $C^o$ expresses $o$ and that $C^o$ is
  stable under it (\ref{aux-consensus-ladder}, stage~1).  So the formalisation
  keeps that bridge and drops the other.
\end{proof}

\aux{Lemma}
\begin{lemma}[Restarting a realisation]
  \label{aux-shift}
  \lean{SocialNetwork.Trajectory.shift, SocialNetwork.Trajectory.state_add,
    SocialNetwork.isGreedyAt_shift}
  \leanok
  \uses{eq2}
  Running a realisation for $m+k$ steps from $u$ is running it for $m$ steps
  and then running the realisation read from step $m$ onwards for $k$ further
  steps; and greediness at step $k$ of the latter is greediness at step $m+k$
  of the former.
\end{lemma}
\begin{proof}
  \leanok
  Induction on $k$.  \textbf{No counterpart in the paper}, which restarts its
  arguments at intermediate times without naming the operation.  Appendix~A
  needs it three times over.
\end{proof}

\section{The probability of a greedy run}

\aux{Lemma}
\begin{lemma}[The gap estimate]
  \label{aux-gap}
  \lean{SocialNetwork.entrySup, SocialNetwork.le_entrySup_sub_one, SocialNetwork.exists_entrySup, SocialNetwork.isGreedyAt_iff_entrySup}
  \leanok
  \uses{eq2}
  Write $y(v) = \max\{v(a,o)\}$.  For $b, o$ with $v(b,o) < y(v)$, the difference
  $v(b,o) - y(v)$ is a strictly negative element of $\pressone\mathbb{Z}$, so
  $v(b,o) - y(v) \le -\pressone$.
\end{lemma}
\begin{proof}
  \leanok
  \textbf{Follows the paper's proof of Proposition~8}, of which this is one
  sentence, extracted so that \ref{prop8} and the biased model can share it.
  Every entry of a state of $\Sspace$ lies in $\mathbb{Z} + \pressone\mathbb{Z}
  = \pressone\mathbb{Z}$.  In the scaled coordinates of \ref{note-scaling} this
  is the statement that two distinct integers differ by at least $1$.
\end{proof}

\paper{Proposition 8}
\begin{proposition}[Probability of a greedy run]
  \label{prop8}
  \lean{SocialNetwork.zeta, SocialNetwork.zeta_le_jumpPMF_argmaxFinset, SocialNetwork.zeta_pow_le_pathMeasure_greedyEvents, SocialNetwork.zeta_le_partialTraj_succ, SocialNetwork.argmaxFinset, SocialNetwork.jumpPMF_apply}
  \leanok
  \uses{aux-greedy, aux-gap, def3}
  For all $m \ge 1$,
  $\mathbb{P}\bigl( \bigcap_{j=1}^{m} \xi_j^u \bigr) \ge (\zeta_\beta)^m$, where
  \[
    \zeta_\beta = \frac{e^{\frac{\beta}{M-1}}}{e^{\frac{\beta}{M-1}} + MN}
      \longrightarrow 1 \quad \text{as } \beta \to +\infty .
  \]
\end{proposition}
\begin{proof}
  \leanok
  \textbf{Follows the paper's proof of Proposition~8} for the bound
  itself.  Two departures, both forced.  The paper's separate treatment of
  $v = 0$ is unnecessary here, the bound below holding uniformly.  And the
  iteration over $m$, which the paper does by conditioning on
  $\Uskel{u}_{m-1} = v$ in~(10), is done as an induction along the
  finite-horizon kernels of the Ionescu--Tulcea construction, Mathlib having no
  decomposition of that shape available; it is the same Markov property, taken
  through the formalism that exists.
  Conditioning on $\Uskel{u}_{m-1} = v$ reduces the claim to a uniform lower
  bound on $\mathbb{P}(\xi_1^v)$.  For $v \ne 0$,
  \[
    \mathbb{P}(\xi_1^v) = \frac{|Y(v)| e^{\beta y(v)}}
      {|Y(v)| e^{\beta y(v)} + \sum_{b,o : v(b,o) < y(v)} e^{\beta v(b,o)}},
  \]
  with $Y(v)$ the set of maximising pairs.  Use $|Y(v)| \ge 1$, \ref{aux-gap}
  for the non-maximising pairs, and $|\actors \times \opinions| = MN$.

  In Lean the iteration is carried out as an induction along the finite-horizon
  kernels of the Ionescu--Tulcea construction: two marginals of one step say
  that the past is almost surely the given history and that the new coordinate
  follows the Gibbs law, and on the intersection of those events a history is
  greedy exactly when its last coordinate maximises the pressure.
\end{proof}

\paper{Remark 4}
\begin{remark}[Bernoulli]
  \label{rem4}
  \lean{SocialNetwork.one_sub_le_zeta_pow, SocialNetwork.one_sub_le_pathMeasure_greedyEvents, SocialNetwork.zeta_pos, SocialNetwork.zeta_le_one, SocialNetwork.one_sub_le_zeta}
  \leanok
  \uses{prop8}
  For any $m \ge 1$ and $\beta \ge 0$,
  $(\zeta_\beta)^m \ge 1 - m M N e^{-\frac{\beta}{M-1}}$.
\end{remark}
\begin{proof}
  \leanok
  $\zeta_\beta = 1/(1+x) \ge 1-x$ with $x = MN e^{-\beta/(M-1)}$.  If $x \ge 1$
  the right-hand side is non-positive and there is nothing to prove; otherwise
  Bernoulli's inequality applies to $(1-x)^m$.
\end{proof}

\section{Concentration of the skeleton measure}

\aux{Lemma}
\begin{lemma}[The greedy run does not revisit $u$]
  \label{aux-greedy-avoids}
  \lean{SocialNetwork.skeleton_ne_of_greedy, SocialNetwork.IsSteepLadder.exists_pos, SocialNetwork.isPositiveAt_of_isSteepLadder, SocialNetwork.isSteepLadder_state_of_greedy}
  \leanok
  \uses{prop7, def4}
  For $u \in \Sspace \setminus \steep$, the greedy run of \ref{prop7} does not
  visit $u$ again: on $\bigcap_{j=1}^{k} \xi_j^u$ with $1 \le k \le (M+1)N$,
  one has $\Uskel{u}_k \ne u$.
\end{lemma}
\begin{proof}
  \leanok
  \uses{prop7, rem5-stable, def4}
  Suppose the run returns, $\Uskel{u}_k = u$ with $k \le (M+1)N$, and repeat its
  first $k$ expressions for ever.  The greedy event at a step reads only the
  matrix and the expression at that step, and the repeated realisation has the
  same state at $n$ as the original at $n \bmod k$ --- which is where the return
  is used --- so the repeated realisation is greedy at \emph{every} step, and it
  sits at $u$ at every multiple of $k$.  By \ref{prop7} it is on $\ladder$ at
  step $(M+1)N$, and a greedy expression on a steep ladder is made from a
  positive entry, so \ref{rem5-stable} keeps it on $\steep$ from there on ---
  including at the multiple $k(M+1)N$.  Hence $u \in \steep$, against the
  hypothesis.

  \textbf{Follows Corollary~8 of [GL24]}, which states this for $M = 2$ and
  proves it ``directly by Part~2 of Proposition~5, by contradiction''.  The
  contradiction is the one above; [GL24] does not write it out, and what it
  needs is the periodicity, which is why the transient is not a formality.
  Greedy runs do return to matrices they have already visited --- from a ladder
  they cycle with period $N$ --- so the hypothesis $u \notin \steep$ is doing
  work, and it does it at the last step.

  \textbf{Supplies a step the paper asserts}, and one it does not: the paper
  reads this off \ref{prop7} --- ``a sequence of events $\xi_j^u$ \dots leads
  the process to $\ladder$, without visiting $u$'' --- and \ref{prop7} says
  where the greedy run ends and nothing about where it passes.  That a greedy
  expression on $\steep$ is a positive expression, which is what lets
  \ref{rem5-stable} apply to a greedy run at all, is
  \texttt{SocialNetwork.IsSteepLadder.exists\_pos}: the $o$-column of a steep
  ladder is non-negative and injective, so with $N \ge 2$ some actor carries
  strictly positive pressure.

  The argument is deterministic, so it is stated on one realisation: no restart
  and no Markov property are needed anywhere in it.
\end{proof}

\paper{Proposition 9}
\begin{proposition}[Concentration of $\muskel$ off the steep ladders]
  \label{prop9}
  \lean{SocialNetwork.measure_le_of_notMem_steepLadderSet, SocialNetwork.kacAvoid, SocialNetwork.kac_identity, SocialNetwork.kac_tsum_le, SocialNetwork.measure_singleton_le_of_avoid, SocialNetwork.kacAvoid_skeletonKernel, SocialNetwork.lintegral_kacAvoid_skeletonKernel, SocialNetwork.returnStep, SocialNetwork.returnBound, SocialNetwork.isSteepLadder_skeleton_of_returnStep, SocialNetwork.returnBound_prod_le, SocialNetwork.one_sub_eta_le, SocialNetwork.eta_lt_one, SocialNetwork.one_div_zeta_le, SocialNetwork.one_add_mul_two_pow_le, SocialNetwork.exp_le_pow_of_one_lt}
  \leanok
  \uses{def4}
  For any $\beta > 0$ and $u \notin \steep$,
  $\muskel(u) \le C' e^{-\beta(N-1)}$ with $C' = (NM)^{(M+1)N+1}$.
\end{proposition}
\begin{proof}
  \leanok
  \uses{prop7, rem5-eta, rem5-eta-iter, prop8, aux-greedy-avoids}
  By Kac's lemma $1/\muskel(u) = \mathbb{E}[\Rskel{u}(u)] \ge \sum_{m \ge
  (M+1)N} \mathbb{P}(\Rskel{u}(u) \ge m)$.  A run of events $\xi_j^u$ takes the
  process to $\ladder$ without visiting $u$ (\ref{prop7} and
  \ref{aux-greedy-avoids}) with the probability of \ref{prop8}, and
  \ref{rem5-eta-iter} then keeps it in $\steep$.  Summing the geometric series
  $\sum_{m \ge 0} \eta^m = 1/(1-\eta)$ and using $1 - \eta \le (1 + MN)
  e^{-\beta(N-1)}$ gives the bound.

  \textbf{Follows the paper's proof}, and departs from it in one place.  The
  paper writes Kac's lemma as the identity $1/\muskel(u) =
  \mathbb{E}[\Rskel{u}(u)]$, which needs the chain to be irreducible and hence
  \ref{thm1-2-skeleton}; the proof uses only
  $\muskel(u)\,\mathbb{E}[\Rskel{u}(u)] \le 1$, and \emph{that} inequality
  holds for every invariant probability measure, with no irreducibility and no
  existence theorem.  It is proved here outright, at the generality of a Markov
  kernel on a countable space (\texttt{SocialNetwork.kac\_tsum\_le}), by
  decomposing the event of visiting $u$ before time $m$ over the last such
  visit.  So this proposition does not wait on Doeblin's criterion, and it is
  stated for an arbitrary invariant probability measure rather than for a named
  $\muskel$.

  \textbf{Supplies two steps the paper asserts.} The composition of \ref{prop7}
  with \ref{rem5-eta-iter} is done on the realisation --- the greedy
  expressions and then the positive ones are one event of the sample space ---
  rather than by restarting the chain at time $(M+1)N$.  And the arithmetic
  that reaches the printed constant needs the regime $MN e^{-\beta/(M-1)} > 1$
  to be treated separately: there $C'$ already exceeds $e^{\beta(N-1)}$ and
  $\muskel(u) \le 1$ suffices.

  \textbf{Adds the hypothesis $u \in \Sspace$}, which the paper works under
  throughout and which \ref{prop7} --- called here --- is stated with.

  \textbf{Rests on} \ref{aux-greedy-avoids}, and through \ref{prop7} on
  \ref{lem19} and~\ref{lem20}.
\end{proof}

\aux{Lemma}
\begin{lemma}[The predecessors of the zero matrix]
  \label{aux-zero-predecessors}
  \lean{SocialNetwork.zeroPredecessor, SocialNetwork.express_zeroPredecessor, SocialNetwork.eq_zeroPredecessor_of_express_eq_zero, SocialNetwork.isState_zeroPredecessor, SocialNetwork.notMem_steepLadderSet_zeroPredecessor, SocialNetwork.eq_of_express_zeroPredecessor_eq_zero, SocialNetwork.skeletonKernel_zeroPredecessor_le, SocialNetwork.measure_exists_zero_row_eq_one}
  \leanok
  \uses{eq2, def3, def4}
  Write $w^{a,o}$ for the matrix whose row $a$ vanishes while every other row
  carries $-1$ on $o$ and $\pressone$ on each of the other opinions.  Then
  \begin{enumerate}
    \item $\muskel$ charges only the matrices with a null row;
    \item a matrix with a null row from which one expression reaches $0$ is one
      of the $NM$ matrices $w^{a,o}$, and $w^{a,o} \in \Sspace \setminus \steep$
      for $N \ge 3$;
    \item the skeleton steps from $w^{a,o}$ to $0$ with probability at most
      $e^{-\frac{\beta}{M-1}}$.
  \end{enumerate}
\end{lemma}
\begin{proof}
  \leanok
  Part 1 is the shape of $\express{a}{o}$ and nothing else: it resets the row of
  the expressing actor, so every matrix the skeleton can reach in one step has a
  null row, and invariance carries that from the kernel to $\muskel$.

  Part 2.  If $\express{a}{o}(v) = 0$ then $v(b,o) = -1$ and $v(b,p) =
  \pressone$ for every $b \ne a$ and $p \ne o$, while the row of $a$ is reset
  and so is left unconstrained --- which is why the null row is needed, and why
  without it the predecessors of $0$ are an infinite family.  With it: no row
  $b \ne a$ vanishes, since $M \ge 2$ makes $-1 \ne 0$, so the null row is the
  row of $a$ and $v = w^{a,o}$.  That $w^{a,o} \in \Sspace$ is the row sum
  $-1 + (M-1)\pressone = 0$.  That $w^{a,o} \notin \steep$ is that for $N \ge 3$
  two actors other than $a$ carry the same pressure on every opinion, where
  \ref{def4} asks the column of the favoured opinion to be injective.

  Part 3.  Read backwards, part 2 says that $(a,o)$ is the only pair whose
  expression carries $w^{a,o}$ to $0$, so the step into $0$ has the probability
  of that one pair.  Its rate is $e^{\beta \cdot 0} = 1$, since the row of $a$
  vanishes, and the normalisation of \ref{def3} already contains the rate
  $e^{\frac{\beta}{M-1}}$ of any actor other than $a$ on any opinion other than
  $o$; the other $NM - 1$ rates in the normalisation are discarded, which is why
  this is a bound and not the probability itself.

  \textbf{Supplies a step the paper asserts.}  \ref{cor10} states parts 2
  and~3 in a single clause and argues neither.  Part 1 has no counterpart at
  all: the paper works throughout in $\Sspace$, where a null row is part of the
  definition, and here it is not assumed but earned --- which is what makes the
  description of the predecessors exhaustive.
\end{proof}

\paper{Corollary 10}
\begin{corollary}[The zero matrix]
  \label{cor10}
  \lean{SocialNetwork.measure_zero_le}
  \leanok
  \uses{prop9}
  $\muskel(0) \le C'' e^{-\beta(N-1+\pressone)}$ with $C'' = (NM) C'$.
\end{corollary}
\begin{proof}
  \leanok
  \uses{prop9, aux-zero-predecessors}
  Invariance writes $\muskel(0)$ as $\sum_v \muskel(v)\,
  \mathbb{P}(\Uskel{v}_1 = 0)$, the mass of each predecessor of $0$ times the
  probability of the step into it.  By
  \ref{aux-zero-predecessors} the sum has $NM$ terms, one per matrix $w^{a,o}$;
  each of them lies in $\Sspace \setminus \steep$ and so carries at most
  $C' e^{-\beta(N-1)}$ by \ref{prop9}; and each step into $0$ costs
  $e^{-\frac{\beta}{M-1}}$.  The three factors give
  $C'' e^{-\beta(N-1+\pressone)}$ with $C'' = (NM) C'$.

  \textbf{Follows the paper's proof}, which reads the extra exponent off the
  maximum $\pressone$ of the states from which $0$ can be entered.  That the
  maximum is what the step costs is the normalisation of the rates of
  \ref{def3}, and that is part 3 of \ref{aux-zero-predecessors}.

  \textbf{Rests on} \ref{prop9}, and through it on \ref{aux-greedy-avoids},
  \ref{lem19} and~\ref{lem20}.
\end{proof}

\paper{Corollary 11}
\begin{corollary}[Hitting a ladder from the zero matrix]
  \label{cor11}
  \lean{SocialNetwork.tendsto_hittingTime_ladderSet_zero}
  \leanok
  \uses{thm2}
  For every fixed $\delta > 0$,
  $\mathbb{P}\bigl(\Rhit{0}(\ladder) > \tau + e^{-\frac{\beta}{M-1}(1-\delta)}\bigr) \to 0$
  as $\beta \to +\infty$, where $\tau$ is exponential of mean $1/(MN)$,
  independent of the process.
\end{corollary}
\begin{proof}
  \textbf{Not formalised}.
\end{proof}

\paper{Remark 6}
\begin{remark}[Reaching $\ladder^o$ before leaving $C^o$]
  \label{rem6}
  \uses{cor11}
  The paper's own note that the first term on the right of equation~(14) may be
  replaced by $\mathbb{P}(\Rhit{u}(\ladder) > t, \Rhit{u}(\ladder^o) <
  \Rhit{u}(C^{-o}))$, whence, by the same steps as part~2 of \ref{thm2}, for
  every fixed $o$ and $\delta > 0$,
  \[
    \inf_{u \in C^o} \mathbb{P}\Bigl(
      \Rhit{u}(\ladder^o) < \min\bigl\{e^{-\frac{\beta}{M-1}(1-\delta)},
        \Rhit{u}(C^{-o})\bigr\}\Bigr) \to 1,
    \qquad \beta \to +\infty .
  \]

  \textbf{Not stated in Lean.} It is the one numbered statement of the paper
  with no counterpart here, and it is recorded so that the count is honest.
  Nothing downstream uses it: it strengthens \ref{cor11} in the direction
  \ref{lem14} needs, and both are blocked on the same continuous-time analysis.
\end{remark}

\section{Metastability}

\paper{Proposition 12}
\begin{proposition}[The exponential-approximation criterion]
  \label{prop12}
  \lean{SocialNetwork.exitTime_approx_exponential, SocialNetwork.IsCharacteristicTime}
  \leanok
  \uses{def2, def1}
  Fix $o$.  Let $c_\beta$ be the positive real with
  $\mathbb{P}(\Rhit{l}(C^{-o}) > c_\beta) = e^{-1}$ for $l \in \ladder^o$.  For
  $\varepsilon_1, \varepsilon_2, s_1, s_2 > 0$ with
  $\varepsilon_1 + \varepsilon_2 \le 1/2$, assume
  \begin{align}
    \varepsilon_1 &\ge \sup_{l \in \ladder^o} \mathbb{P}(\Rhit{l}(C^{-o}) \le s_1),
      \label{eq:h15} \\
    \varepsilon_2 &\ge \sup_{u \in \Sspace} \mathbb{P}(\Rhit{u}(\ladder^o \cup C^{-o}) > s_2),
      \label{eq:h16} \\
    (s_2/c_\beta) \vee \varepsilon_2 &\le C e^{-\delta\beta}, \label{eq:h17} \\
    \sup_{u \in C^o} \mathbb{P}(\Rhit{u}(C^{-o}) \le s_2) &\le K e^{-\theta\beta}.
      \label{eq:h18}
  \end{align}
  Then there is $\beta_0(C,\delta)$ and $K'(\delta,C,K,\theta)$ such that for
  $\beta \ge \beta_0$ and $u \in C^o$ the rescaled exit time is exponential up
  to $K' \beta^3 e^{-\min(\delta/3, 1/2, \theta)\beta}$, and the mean exit times
  from two states of $C^o$ agree to the same order.
\end{proposition}
\begin{proof}
  This is a consequence of Theorem~5.3 of [LM22].  \textbf{Declared as an
  axiom}, not proved.  It is not a result of this paper: its content is a
  metastability estimate for a general time-homogeneous strong Markov
  process, and nothing inside this library can discharge it.  Making it an
  \texttt{axiom} rather than a \texttt{sorry} states that plainly, and the CI
  fails if any result claimed complete reaches it.

  The Lean statement is not the paper's display verbatim.  In the paper
  $\varepsilon_1, \varepsilon_2, s_1, s_2$ are introduced before $\beta$, but
  Section~5.3 instantiates them at $s_2 = 2\beta$ and
  $\varepsilon_2 = (M+1)^2N^2 e^{-\beta/((M+1)N)}$, and notes that
  $\varepsilon_1 + \varepsilon_2 < 1/2$ holds only ``for $\beta$ sufficiently
  big''.  Bound as constants ahead of $\beta$ the hypotheses are
  \emph{unsatisfiable}, so in Lean they are functions of $\beta$ and
  \eqref{eq:h15}--\eqref{eq:h18} are required only above a threshold
  $\beta_1$.

  The axiom is also attached to \emph{this} process on purpose.  Stated for
  an arbitrary family of measures and an arbitrary hitting time it would be
  \textbf{inconsistent}: the zero measure with $\ladder^o = \emptyset$
  satisfies \eqref{eq:h15}--\eqref{eq:h18} vacuously and falsifies the
  conclusion at $t = 0$.  What excludes that is the strong Markov property,
  which is the content of [LM22] and is not expressible here.  Theorem~31
  therefore needs its own twin; one axiom cannot serve both models.
\end{proof}

\aux{Lemma}
\begin{lemma}[Inside the proof of part~2 of Theorem~2]
  \label{aux-hitting-rate}
  \lean{SocialNetwork.probHittingGT_ladderSet_le_of_ne_zero}
  \leanok
  \uses{def1}
  For any $u \in \Sspace \setminus \{0\}$ and any $t > 0$,
  \[
    \mathbb{P}(\Rhit{u}(\ladder) > t)
      \;\le\; 1 - \zeta_\beta^{(M+1)N}
        + (M+1)N \exp\!\left(-\frac{e^{\beta/(M-1)}\, t}{(M+1)N}\right).
  \]
\end{lemma}
\begin{proof}
  \uses{prop6, prop8}
  \textbf{Not formalised, and not a numbered statement of the paper}: this
  inequality is displayed inside the proof of part~2 of \ref{thm2}, and
  Theorem~2.2 is the limit it gives at $t = e^{-\beta(1-\delta)/(M-1)}$.  It is
  stated here because \ref{lem13} uses the inequality and not the limit: taking
  the limit throws the rate away, so Lemma~13 cannot rest on Theorem~2.2 as
  stated.  See Section~\ref{sec:corrections}.
\end{proof}

\paper{Equation (19)}
\begin{lemma}[Displayed inside the proof of Lemma~13]
  \label{eq19}
  \lean{SocialNetwork.probHittingGT_ladderSet_zero_le}
  \leanok
  \uses{def1}
  \[
    \mathbb{P}(\Rhit{0}(\ladder) > 2\beta)
      \;\le\; \mathbb{P}(\tau > \beta)
        + \sup_{u \ne 0} \mathbb{P}(\Rhit{u}(\ladder) > \beta),
  \]
  with $\mathbb{P}(\tau > \beta)$ written as the paper evaluates it,
  $e^{-\beta/(MN)}$.
\end{lemma}
\begin{proof}
  \uses{cor11}
  \textbf{Not formalised, and not a numbered statement of the paper}: it is
  displayed as equation~(19) inside the proof of \ref{lem13} and attributed
  there to \ref{cor11}, which is itself stated as a limit.  Note that $\tau$ is
  declared exponential of mean $1/(MN)$, for which $\mathbb{P}(\tau > \beta) =
  e^{-MN\beta}$ and not $e^{-\beta/(MN)}$; since $e^{-MN\beta} \le
  e^{-\beta/(MN)}$ for $\beta \ge 0$, the form written here is the weaker of
  the two and Lemma~13 follows from either reading.
\end{proof}

\paper{Lemma 13}
\begin{lemma}[Hitting a ladder within $2\beta$]
  \label{lem13}
  \lean{SocialNetwork.probHittingGT_ladderSet_le}
  \leanok
  \uses{def1}
  For any $u \in \Sspace$,
  $\mathbb{P}(\Rhit{u}(\ladder) > 2\beta) \le (M+1)^2 N^2 e^{-\beta/((M+1)N)}$.
  This is \eqref{eq:h16} with $s_2 = 2\beta$.
\end{lemma}
\begin{proof}
  \leanok
  \uses{eq19, aux-hitting-rate, rem4}
  Split according to whether the start is $0$.  From $0$, equation~(19) of
  \ref{eq19}; from any other state the event at $2\beta$ is contained in the
  event at $\beta$.  Both are then bounded by \ref{aux-hitting-rate} at $t =
  \beta$, and $1 - \zeta_\beta^{(M+1)N}$ by \ref{rem4}. \emph{Formalised, but
  not sorry-free}: the proof is complete in Lean and inherits \texttt{sorryAx}
  from the two displays above alone. \textbf{Supplies a step the paper
  asserts}: ``putting the inequalities above together'' is the comparison of
  each of the three terms with $e^{-\beta/((M+1)N)}$ --- using $MN \le (M+1)N$,
  $M-1 \le (M+1)N$ and $e^{\beta/(M-1)} \ge 1$ --- which leaves the integer
  inequality $1 + (M+1)N \le (M+1)N^2$, true since $N \ge 3$.
\end{proof}

\paper{Lemma 14}
\begin{lemma}[Two-sided control of the exit from a consensus set]
  \label{lem14}
  \lean{SocialNetwork.le_probHittingGT_consensusOther, SocialNetwork.probHittingLE_consensusOther_le}
  \leanok
  \uses{def2, def1}
  For any $t > 0$ and any $o$:
  \begin{enumerate}
    \item for $l \in \ladder^o$,
      $\mathbb{P}(\Rhit{l}(C^{-o}) > t) \ge e^{-2tN^3(M+1)^3 e^{-\beta/(M-1)}}$;
    \item for $u \in C^o$,
      $\mathbb{P}(\Rhit{u}(C^{-o}) \le t) \le (N^2M + 2tN^3(M+1)^3) e^{-\beta/(M-1)}$.
  \end{enumerate}
\end{lemma}
\begin{proof}
  Proved in Appendix~B of the paper.  \textbf{Not formalised}.
\end{proof}

\paper{Corollary 15}
\begin{corollary}[A lower bound for $c_\beta$]
  \label{cor15}
  \lean{SocialNetwork.le_characteristicTime}
  \leanok
  \uses{prop12, lem14}
  For any $\beta \ge 0$,
  $c_\beta \ge \tfrac{1}{2} N^{-3}(M+1)^{-3} e^{\frac{\beta}{M-1}}$.
\end{corollary}
\begin{proof}
  \leanok
  \uses{lem14, aux-ladder-nonempty}
  Apply \ref{lem14}, part 1, at $t = c_\beta$ to the ladder $l$ of
  \ref{aux-ladder-nonempty}, so that
  \[
    e^{-1} \;=\; \mathbb{P}(\Rhit{l}(C^{-o}) > c_\beta)
    \;\ge\; e^{-2c_\beta N^3(M+1)^3 e^{-\beta/(M-1)}} ,
  \]
  and rearrange. \emph{Formalised, but not sorry-free}: the proof above is
  complete in Lean, and inherits \texttt{sorryAx} from \ref{lem14} alone.  Its
  only addition to the paper is the witness: the statement quantifies over $l
  \in L^o$, so it is vacuous unless $L^o \ne \emptyset$, which the paper does
  not record.
\end{proof}

\aux{Lemma}
\begin{lemma}[$\ladder^o$ is inhabited]
  \label{aux-ladder-nonempty}
  \lean{SocialNetwork.ladderOf, SocialNetwork.isLadder_ladderOf,
    SocialNetwork.exists_isLadder, SocialNetwork.ladderSet_nonempty}
  \leanok
  \uses{def1}
  For every $o \in O$ the set $L^o$ is non-empty: the matrix
  $l(a, o) = a$, $l(a, p) = -a/(M-1)$ for $p \ne o$, is a ladder supporting
  $o$.
\end{lemma}
\begin{proof}
  \leanok
  Each row sums to $a - (M-1)\tfrac{a}{M-1} = 0$, actor $1$ carries the null
  row, the column for $o$ is the staircase $\{0, \dots, N-1\}$ by
  construction, and the other columns are the prescribed multiples.
  \textbf{Supplies a step the paper asserts}: every statement of the form
  ``for any $l \in L^o$'' concluding something about $L^o$ --- Corollary~15
  is the first --- is vacuous without it.
\end{proof}

\section{The negative-bias regime}

\paper{Theorem 16}
\begin{theorem}[The biased process does not explode]
  \label{thm16}
  \lean{SocialNetwork.Bias.biasedNonExplosion}
  \leanok
  \uses{eq7}
  For $\beta \ge 0$, $\alpha < 0$ and $u \in \Salpha$, the biased process does
  not explode.
\end{theorem}
\begin{proof}
  \uses{prop21, thm1-1}
  As \ref{thm1-1}, with \ref{prop21} in place of \ref{prop5}. \textbf{Not
  formalised}.
\end{proof}

\paper{Proposition 17}
\begin{proposition}[Reaching bounded pressures]
  \label{prop17}
  \lean{SocialNetwork.Bias.measure_biasedBounded_ge, SocialNetwork.Bias.biasedBounded}
  \leanok
  \uses{eq7}
  For $\alpha < 0$, $\beta \ge 0$ and $u \in \Salpha$,
  $\mathbb{P}(U^{\alpha,u}_{T_N} \in B_N^\alpha) \ge (NM)^{-N}$, where
  $B_N^\alpha = \{ u \in \Salpha : \max_{a,o} u(a,o) \le N \}$.
\end{proposition}
\begin{proof}
  \leanok
  \uses{prop6, prop8}
  As \ref{prop6}, $\bigcap_{j=1}^N \xi_j^{\alpha,u} \subset
  \{U^{\alpha,u}_{T_N} \in B_N^\alpha\}$, and each $\xi_j^{\alpha,u}$ is the
  most probable of $NM$ choices. \textbf{Follows the paper's proof of
  Proposition~17}, in its two halves.

  The inclusion is \ref{prop6}'s argument with $n_a$ in place of
  $\|U(a,\cdot)\|_\infty$, in its two cases: either the first $N$ expressions
  come from $N$ distinct actors, and every row was reset, or some actor
  expresses at $j < k < N$ and greediness at $k$ caps the whole profile at $n_a
  = k-j-1$, which $N-k$ further expressions raise to at most $N-j-1$.  Unlike
  \ref{prop22}, the chain closes here: greediness is \emph{exact}, so the
  maximum at the repeat time is the expressing actor's own entry, with no slack
  of $1/(2\gamma)$ to absorb.

  The probability is the one-step bound: with $\beta \ge 0$ the maximising
  pairs carry the largest rate, so they carry at least $(NM)^{-1}$ of the total
  weight, and the iteration is the one of \ref{prop8}.  That iteration is run
  once in Lean, for an arbitrary choice of admissible pairs at each profile,
  and instantiated both here and at \ref{prop24}.
\end{proof}

\paper{Proposition 18}
\begin{proposition}[One actor expresses for ever]
  \label{prop18}
  \lean{SocialNetwork.Bias.inf_measure_forall_eq_first_pos,
    SocialNetwork.Bias.exists_pos_le_pathMeasure_forall_lt_fst,
    SocialNetwork.Bias.soloPath,
    SocialNetwork.Bias.pressure_stateAfter_soloPath_of_ne,
    SocialNetwork.Bias.soloOtherRate,
    SocialNetwork.Bias.inv_le_biasedJumpPMF_soloPath,
    SocialNetwork.Bias.prod_le_pathMeasure_cylinder,
    SocialNetwork.Bias.prod_inv_add_ge,
    SocialNetwork.Bias.sum_le_of_le_geometric,
    SocialNetwork.occCount,
    SocialNetwork.uniformSeq,
    SocialNetwork.freqGood,
    SocialNetwork.tendsto_occCount_ae,
    SocialNetwork.exists_uniformSeq_freqGood_pos,
    SocialNetwork.uniformSeq_freqGood_le}
  \leanok
  \uses{prop17}
  For $\alpha < 0$ and $\beta > 0$,
  $\inf_{u \in B_N^\alpha} \mathbb{P}\bigl( \bigcap_{j \ge 1} \{A_j^u = A_1^u\} \bigr) > 0$.
\end{proposition}
\begin{proof}
  \leanok
  \textbf{Follows the paper's proof of Proposition~18.}

  Fix $u \in B_N^\alpha$ and an actor $a$ whose row is null.  Equation~(6)
  provides one --- it is the condition $\min_a \|u(a,\cdot)\|_\infty = 0$ ---
  so the paper's ``without loss of generality $u(1,o) = 0$ for all $o$'' is a
  choice of actor and nothing more.

  Along the run in which $a$ expresses the word $o_1, o_2, \dots$ and no other
  actor ever expresses, the memory of $a$ is reset at every step and every
  other row is the initial one shifted by the counts of that word.  So the pair
  $(a, o_m)$ carries rate $1$ while the whole profile carries at most
  $M + \lambda_m$, with the $\lambda_m$ of~(22).  Equation~(21) is then the
  product of these one-step probabilities, summed over the $M^n$ words of
  length $n$: the measure of one cylinder is at least the product along it, and
  the cylinders of distinct words are disjoint.

  Since $Me^{\lambda} \ge M + \lambda$ --- the paper's $\ln(1+x) \ge x/(1+x)$
  at $x = -\lambda/(M+\lambda)$, composed with its own weakening of
  $e^{-\lambda/M}$ to $e^{-\lambda}$ --- the product along a word is at least
  $M^{-n} e^{-\sum_{m < n} \lambda_m}$.  Restricting the sum to the words that
  satisfy the constraint of $E_\varepsilon^k$ up to time $n$ makes
  $\sum_{m<n} \lambda_m$ bounded by a constant, uniformly in $n$: above $k$ the
  exponent of~(22) is at most $\beta(N - c(\alpha,\varepsilon)\,m)$, which
  is~(24), and $c(\alpha,\varepsilon) > 0$ for $\varepsilon$ small is exactly where
  $\alpha < 0$ is used.  Those words carry at least the uniform mass
  $\mathbb{P}(E_\varepsilon^k)$, which the strong law of large numbers makes
  positive for some $k$.  The passage from $\bigcap_{m \le n}$ to
  $\bigcap_{m \ge 1}$ is continuity from above along decreasing events.

  \textbf{Supplies two steps the paper asserts}: the choice of actor above, and
  the terms of index below $k$, which $E_\varepsilon^k$ does not constrain and
  which are bounded crudely by
  $\lambda_m \le (N-1)\,M\,e^{\beta(N + k(1+\gamma))}$.
\end{proof}

\chapter{Appendix A: the road to Proposition~7}

\paper{Definition 5}
\begin{definition}[Matrices favouring an opinion]
  \label{def5-favouring}
  \lean{SocialNetwork.IsFavouring, SocialNetwork.IsFavouringWith, SocialNetwork.favouringSet}
  \leanok
  \uses{eq2}
  $u \in S^o$ when there are $r \in [0,1)$, $n(u) \in \{1, \dots, N-1\}$ and
  $n(u)$ distinct actors $a_1(u), \dots, a_{n(u)}(u)$ with
  \[
    u(a_j(u), o) \ge j + r \quad \forall j, \qquad
    u(a,p) \le n(u) + r - \pressone \quad \forall a, \forall p \ne o .
  \]
  In scaled coordinates $r$ becomes an integer $\rho \in [0, M-1)$, which is
  exactly $r \in [0,1) \cap \pressone\mathbb{Z}$ — the only case the paper's own
  construction produces, since it takes $r = m - \lfloor m \rfloor$ for $m$ an
  entry.
\end{definition}

\paper{Definition 5}
\begin{definition}[The first repeat time]
  \label{def5-repeat}
  \lean{SocialNetwork.firstRepeat, SocialNetwork.firstRepeat_le, SocialNetwork.actor_injOn_lt_firstRepeat, SocialNetwork.exists_repeat}
  \leanok
  \uses{aux-trajectory}
  $\tau(u) := \inf\{ n \ge 1 : A_n \in \{A_1, \dots, A_{n-1}\} \}$, the number of
  expressions until an actor expresses for the second time.  Then
  $\tau(u) \in \{2, \dots, N+1\}$.
\end{definition}
\begin{proof}
  \leanok
  Pigeonhole: among $A_1, \dots, A_{N+1}$ two must coincide.
\end{proof}

\paper{Lemma 19}
\begin{lemma}[The first repeat time reaches $\bigcup_o S^o$]
  \label{lem19}
  \lean{SocialNetwork.isFavouring_state_firstRepeat}
  \leanok
  \uses{def5-favouring, def5-repeat, aux-greedy}
  For any $u \in \Sspace$, the event $\xi^u_{\tau(u)}$ implies
  $\Uskel{u}_{\tau(u)} \in \bigcup_{o} S^o$.
\end{lemma}
\begin{proof}
  \textbf{Not formalised, and not blocked on Mathlib}: the written proof
  asserts a sequence of $\lfloor m \rfloor + 1$ distinct actors ``by (25)''
  without giving the construction, and rules out $m = 0$ through $\tau(u) = 2$
  rather than through $m = 0$ itself.  See \ref{note-lem19} for a construction
  that works and the corrected case split.
\end{proof}

\paper{Lemma 20}
\begin{lemma}[From a favouring matrix to a consensus state]
  \label{lem20}
  \lean{SocialNetwork.isConsensus_state_of_favouring}
  \leanok
  \uses{def5-favouring, def2, aux-greedy}
  For any $u \in S^o$, the event $\bigcap_{j=1}^{(M-1)(n(u)+1)-1} \xi_j^u$
  implies $\Uskel{u}_{(M-1)(n(u)+1)-1} \in C^o$.
\end{lemma}
\begin{proof}
  \uses{lem20-opening, lem20-closing}
  The proof opens with \ref{lem20-opening} and closes with \ref{lem20-closing};
  its induction is the part that fails, see \ref{note-lem20}. \textbf{Not
  formalised, and not blocked on Mathlib}.
\end{proof}

\paper{Lemma 20}
\begin{lemma}[The opening step]
  \label{lem20-opening}
  \lean{SocialNetwork.opinion_eq_of_isMax_of_favouring}
  \leanok
  \uses{def5-favouring}
  In a state of $S^o$ every maximising entry lies in column $o$, so a greedy
  expression expresses $o$.
\end{lemma}
\begin{proof}
  \leanok
  \textbf{Supplies a step the paper asserts}: the proof of \ref{lem20} opens
  ``then $\Uskel{u}_0(A_1,O_1) = \max u(a,o)$ and $O_1 = o$'', asserting $O_1 =
  o$.  The witness $a_{n(u)}$ carries at least $n(u) + r$ in column $o$, which
  strictly exceeds the bound $n(u) + r - \pressone$ imposed on every other
  column.
\end{proof}

\paper{Lemma 20}
\begin{lemma}[The closing step]
  \label{lem20-closing}
  \lean{SocialNetwork.isConsensus_of_nonpos}
  \leanok
  \uses{def2, eq2}
  A non-zero state of $\Sspace$ whose columns other than $o$ are non-positive is
  a consensus state for $o$.
\end{lemma}
\begin{proof}
  \leanok
  \textbf{Supplies a step the paper asserts}: the proof of \ref{lem20} ends
  ``by taking $k = (M-1)(n(u)+1)-1$ we conclude the proof'', the passage to
  $C^o$ being left to the reader.  The vanishing row sums turn the sign
  condition on the other columns into non-negativity on column $o$.
\end{proof}

\chapter{Appendix C: the biased analogues}

\paper{Proposition 21}
\begin{proposition}[Low pressure recurs, under bias]
  \label{prop21}
  \lean{SocialNetwork.Bias.exists_pressure_lt}
  \leanok
  \uses{eq6}
  For any $u \in \Salpha$,
  $\inf\{ n \ge 1 : \|U^{\alpha,u}_{T_{n-1}}(A_n^\alpha, \cdot)\|_\infty < N \} \le N$.
\end{proposition}
\begin{proof}
  \leanok
  \uses{prop5}
  The proof of \ref{prop5} is combinatorial and does not see the bias.
  \textbf{Follows the paper's proof of Proposition~5}, which Section~5.4
  invokes verbatim: if each of the first $N$ expressions came from an actor
  carrying pressure at least $N$, then none of them is the actor whose row is
  null, and no two of them coincide, exhibiting $N+1$ distinct actors in a
  network of $N$.

  One step is supplied.  The quantity that is reset by expressing and then grows
  by at most one per step is, in the unbiased model,
  $\|U(a, \cdot)\|_\infty$; here it is $n_a$, the number of expressions heard.
  The passage between them is $u(a,p) = c_p(1+\gamma) - \gamma n_a \le n_a$,
  since $c_p \le n_a$ --- immediate from equation~(6), but not recorded there.
\end{proof}

\paper{Remark 7}
\begin{definition}[The near-greedy event]
  \label{rem7}
  \lean{SocialNetwork.Bias.IsNearGreedyAt, SocialNetwork.Bias.nearGreedyEvents, SocialNetwork.Bias.IsBiasedGreedyAt, SocialNetwork.Bias.biasedGreedyEvents, SocialNetwork.Bias.ofHistoryPath, SocialNetwork.Bias.stateAfter_ofHistoryPath_frestrictLe, SocialNetwork.Bias.measurableSet_nearGreedyEvents, SocialNetwork.Bias.measurableSet_biasedGreedyEvents}
  \leanok
  \uses{eq6}
  \[
    \tilde{\xi}_n^{\alpha,u} = \Bigl\{ U^{\alpha,u}_{T_{n-1}}(A_n^\alpha, O_n^\alpha)
      > \max_{(a,o)} U^{\alpha,u}_{T_{n-1}}(a,o) - \tfrac{1}{2\gamma} \Bigr\} .
  \]
  The slack of $\tfrac{1}{2\gamma}$ is what replaces the lattice gap $\pressone$
  of the unbiased model: it is what makes a uniform lower bound on
  $\mathbb{P}(\tilde{\xi}_n^{\alpha,u})$ available.

  The two measurability lemmas named above have \textbf{no counterpart in the
  paper} (\ref{note-green}).
\end{definition}

\paper{Proposition 22}
\begin{proposition}[Confinement under bias]
  \label{prop22}
  \lean{SocialNetwork.Bias.entry_mem_of_nearGreedy}
  \leanok
  \uses{rem7}
  For any $u \in \Salpha$,
  $\bigcap_{j=1}^{N} \tilde{\xi}_j^{\alpha,u} \subset \{ -MN < U^{\alpha,u}_{T_N} < N \}$.
\end{proposition}
\begin{proof}
  \uses{prop6}
  \textbf{Not formalised, and it does not follow from Proposition~6 as the text
  asserts.} See \ref{note-prop22}.
\end{proof}

\paper{Proposition 23}
\begin{proposition}[Reaching a biased ladder]
  \label{prop23}
  \lean{SocialNetwork.Bias.exists_horizon_isBiasedLadder}
  \leanok
  \uses{rem7, eq8}
  There is $C(\alpha, M, N) \in \mathbb{N}$ with
  $\bigcap_{j=1}^{C} \tilde{\xi}_j^{\alpha,u}
   \subset \{ U^{\alpha,u}_{T_{C}} \in \ladder_\alpha \}$.
\end{proposition}
\begin{proof}
  \uses{prop7}
  \textbf{Not formalised}.
\end{proof}

\paper{Proposition 24}
\begin{proposition}[Probability of a biased greedy run]
  \label{prop24}
  \lean{SocialNetwork.Bias.biasedZeta, SocialNetwork.Bias.biasedZeta_pow_le}
  \leanok
  \uses{rem7}
  For all $m \ge 1$,
  $\mathbb{P}\bigl( \bigcap_{j=1}^{m} \tilde{\xi}_j^{\alpha,u} \bigr)
   \ge (\zeta_{\alpha,\beta})^m$ with
  $\zeta_{\alpha,\beta} = e^{\frac{\beta}{2\gamma}} /
   \bigl( e^{\frac{\beta}{2\gamma}} + MN \bigr)$.
\end{proposition}
\begin{proof}
  \leanok
  \uses{prop8}
  As \ref{prop8}, with the lattice gap replaced by the slack of \ref{rem7}.
  \textbf{Follows the paper's proof of Proposition~8}, which Appendix~C invokes
  for this statement.

  The one-step bound is the same computation: the maximising pair carries
  weight $e^{\beta y}$, every pair outside $\tilde{Y}_\gamma$ carries at most
  $e^{\beta y} e^{-\beta/(2\gamma)}$, and there are at most $MN$ of them.  What
  the lattice supplied in the unbiased model --- a gap of $1/(M-1)$ below the
  maximum --- the event $\tilde{\xi}$ supplies by construction, which is what
  Remark~7 is for.  As in \ref{prop8} the iteration goes through the
  finite-horizon kernels of the Ionescu--Tulcea construction rather than the
  conditioning of~(10).
\end{proof}

\paper{Theorem 25}
\begin{theorem}[Existence and uniqueness of $\mu_{\beta,\alpha}$]
  \label{thm25}
  \lean{SocialNetwork.Bias.existsUnique_biasedInvariant}
  \leanok
  \uses{eq7}
  For $\beta > 0$ and $0 < \alpha < \pressone$, the biased process does not
  explode and has a unique invariant probability measure $\mu_{\beta,\alpha}$.
\end{theorem}
\begin{proof}
  \uses{thm1-2}
  As \ref{thm1-2}.  \textbf{Not formalised}.
\end{proof}

\paper{Proposition 26}
\begin{proposition}[Concentration off the biased steep ladders]
  \label{prop26}
  \lean{SocialNetwork.Bias.biasedMeasure_le_of_notMem_steepLadder}
  \leanok
  \uses{eq8, thm25}
  For $0 < \alpha < \pressone$, $\beta > 0$ and $u \notin \steep_\alpha$,
  $\tilde{\mu}_{\alpha,\beta}(u) \le \tilde{C} e^{-\beta(N-1)}$.
\end{proposition}
\begin{proof}
  \uses{prop9, rem8}
  As \ref{prop9}, using \ref{rem8} and $0 \notin \Salpha$. \textbf{Not
  formalised}.
\end{proof}

\paper{Theorem 27}
\begin{theorem}[Concentration of $\mu_{\alpha,\beta}$]
  \label{thm27}
  \lean{SocialNetwork.Bias.biasedMeasure_ladderSet_ge, SocialNetwork.Bias.tendsto_biasedHittingTime}
  \leanok
  \uses{eq8}
  For $0 < \alpha < \pressone$ there is $C > 0$ with
  $\mu_{\alpha,\beta}(\ladder_\alpha) \ge 1 - C e^{-\beta(M-1)\alpha}$, and for
  every fixed $\delta > 0$,
  $\sup_{u \in \Salpha} \mathbb{P}\bigl( R^{\alpha,\beta,u}(\ladder_\alpha)
   > e^{-\beta(M-1)\alpha(1-\delta)} \bigr) \to 0$.
\end{theorem}
\begin{proof}
  \uses{thm2, prop26, rem8}
  As \ref{thm2}: the lower bound for the jump rate of a non-null state is
  $e^{\beta(M-1)\alpha}$ by \ref{rem8}, and the contribution of the zero matrix
  vanishes because $0 \notin \Salpha$. \textbf{Not formalised}.
\end{proof}

\paper{Lemma 28}
\begin{lemma}[Hitting a biased ladder within $2\beta$]
  \label{lem28}
  \lean{SocialNetwork.Bias.biasedProbHitting_le}
  \leanok
  \uses{eq8}
  For any $u \in \Salpha$,
  $\mathbb{P}(R^{\alpha,\beta,u}(\ladder_\alpha) > 2\beta) \le C e^{-\frac{\beta}{2\gamma}}$,
  with $C$ depending only on $\alpha$, $M$ and $N$.
\end{lemma}
\begin{proof}
  \uses{lem13}
  \textbf{Not formalised.} \emph{Restated}: the Lean statement of this lemma
  used to be about the skeleton path measure and the discrete steps $k \le
  \lceil 2\beta \rceil$ rather than about the continuous-time hitting time
  $R^{\alpha,\beta,u}$, and it bound $C$ after $\beta$ and $u$, so the constant
  was free to depend on both.  Neither matches the display above, and neither
  can serve as \eqref{eq:h16} for \ref{prop12-biased}, which is what this lemma
  exists for.  It now has the shape of \ref{lem13} with $\pressone$ replaced by
  $\tfrac{1}{2\gamma}$.  A formalisation-side correction, not a correction to
  the paper.
\end{proof}

\paper{Lemma 29}
\begin{lemma}[Two-sided control of the exit, under bias]
  \label{lem29}
  \lean{SocialNetwork.Bias.le_biasedProbHittingGT, SocialNetwork.Bias.biasedProbHittingLE_le}
  \leanok
  \uses{eq8}
  The two-sided control of \ref{lem14}, with $\pressone$ replaced by
  $\tfrac{1}{2\gamma}$.
\end{lemma}
\begin{proof}
  \uses{lem14}
  \textbf{Not formalised}.
\end{proof}

\paper{Corollary 30}
\begin{corollary}[A lower bound for $c_{\alpha,\beta}$]
  \label{cor30}
  \lean{SocialNetwork.Bias.le_biasedCharacteristicTime, SocialNetwork.Bias.IsBiasedCharacteristicTime}
  \leanok
  \uses{lem29}
  For $0 < \alpha < \pressone$ and $\beta \ge 0$,
  $c_{\alpha,\beta} \ge \tfrac{1}{2} N^{-3}(M+1)^{-3} e^{\frac{\beta}{2\gamma}}$.
\end{corollary}
\begin{proof}
  \leanok
  \uses{cor15, lem29, aux-biased-ladder-nonempty}
  The rearrangement of \ref{cor15}, with $\pressone$ replaced by
  $\tfrac{1}{2\gamma}$: apply \ref{lem29}, part~1, at $t = c_{\alpha,\beta}$ to
  the ladder of \ref{aux-biased-ladder-nonempty}. \emph{Formalised, but not
  sorry-free}: complete in Lean, inheriting \texttt{sorryAx} from \ref{lem29}
  alone.  Its only addition to the paper is the witness \ref{cor15} needed as
  well.
\end{proof}

\aux{Lemma}
\begin{lemma}[$\ladder_\alpha^o$ is inhabited]
  \label{aux-biased-ladder-nonempty}
  \lean{SocialNetwork.Bias.biasedLadderOf,
    SocialNetwork.Bias.pressure_biasedLadderOf,
    SocialNetwork.Bias.isBiasedLadder_biasedLadderOf,
    SocialNetwork.Bias.exists_isBiasedLadder,
    SocialNetwork.Bias.biasedLadderSet_nonempty}
  \leanok
  \uses{eq8}
  For every $o \in O$ the set $L_\alpha^o$ is non-empty: the profile in which
  actor $a$ has heard exactly $a$ expressions, all of them of $o$, is a biased
  ladder supporting $o$.
\end{lemma}
\begin{proof}
  \leanok
  That memory has $c_o = n_a = a$ and $c_p = 0$ for $p \ne o$, so $u(a,o) =
  a(1+\gamma) - \gamma a = a$ and $u(a,p) = -\gamma a$, which is equation~(9).
  Actor $1$ has heard nothing and the $n_a$ are pairwise distinct, so the
  profile lies in $S^\alpha$. \textbf{No counterpart in the paper}, exactly as
  for \ref{aux-ladder-nonempty}: without it \ref{cor30} is vacuous.
\end{proof}

\paper{Theorem 31}
\begin{theorem}[Metastability under bias]
  \label{thm31}
  \lean{SocialNetwork.Bias.biasedMetastability}
  \leanok
  \uses{eq8}
  For $0 < \alpha < \pressone$ there are $\beta_0, C_1 > 0$ and $C_2 > 0$,
  depending only on $\alpha$, $M$ and $N$, such that the rescaled exit time from
  a biased consensus set is exponential of parameter one up to
  $C_1 \beta^3 e^{-C_2\beta}$.
\end{theorem}
\begin{proof}
  \leanok
  \uses{prop12-biased, thm3, cor30,
    lem28, lem29, aux-biased-ladder-consensus}
  As \ref{thm3}, which is all Appendix~C says. \emph{Formalised, but neither
  sorry-free nor axiom-free}: the assembly is complete in Lean, and inherits
  \texttt{sorryAx} from \ref{lem28} and~\ref{lem29} together with the second
  [LM22] axiom, the biased twin \ref{prop12-biased}.

  The four assumptions are checked as in Section~5.3, with $\pressone$ replaced
  by $\tfrac{1}{2\gamma}$ throughout:
  \begin{itemize}
    \item \eqref{eq:h15} at $s_1 = 1$,
      $\varepsilon_1 = 2N^3(M+1)^3 e^{-\beta/(2\gamma)}$, from part 1 of
      \ref{lem29} through $1 - e^{-x} \le x$;
    \item \eqref{eq:h16} at $s_2 = 2\beta$, from
      \ref{lem28} and the inclusion $\ladder_\alpha \subseteq \ladder_\alpha^o
      \cup C_\alpha^{-o}$ of \ref{aux-biased-ladder-consensus};
    \item \eqref{eq:h17} at $\delta = 1/(4\gamma)$, from
      \ref{cor30};
    \item \eqref{eq:h18} at $\theta = 1/(4\gamma)$, from part 2 of
      \ref{lem29} at $t = 2\beta$.
  \end{itemize}

  \textbf{Supplies a step the paper asserts}: Appendix~C names none of the
  constants.  They are what Section~5.3 produces once the exponent
  $\tfrac{1}{2\gamma}$ replaces $\pressone$ --- halved to $\tfrac{1}{4\gamma}$
  to absorb the factor $\beta$ of step~(20), so that $\delta = \theta$ here
  where the unbiased proof had two different values --- together with an
  explicit threshold $\beta_1 = \max\{1, 4\gamma(2N^3(M+1)^3 + C)\}$ above which
  $\varepsilon_1 + \varepsilon_2 \le 1/2$, and, as in Theorem~3, constants taken
  uniform over the finitely many opinions.
\end{proof}

\paper{Proposition 12}
\begin{proposition}[For the biased process]
  \label{prop12-biased}
  \lean{SocialNetwork.Bias.biasedExitTime_approx_exponential}
  \leanok
  \uses{eq8}
  Under \eqref{eq:h15}--\eqref{eq:h18} read for the biased process, the
  rescaled exit time from $C_\alpha^o$ is exponential of parameter one up to
  $K'\beta^3 e^{-\min(\delta/3,\,1/2,\,\theta)\beta}$, and the mean exit time
  barely depends on the starting profile.
\end{proposition}
\begin{proof}
  \uses{prop12}
  \textbf{An axiom, not a theorem, and the second one this development asks the
  reader to trust.} Appendix~C never states it: it says only that ``the proof
  of Theorem~31 follows exactly as the proof of Theorem~3'', and that proof
  runs through \ref{prop12}, which is stated for the unbiased process alone.

  One axiom cannot serve both models, for the reason recorded at \ref{prop12}:
  stated abstractly the assertion is \textbf{inconsistent}, so it has to be
  attached to a concrete process, and the unbiased axiom is attached to the
  unbiased one.  Everything resting on this one is listed separately in the CI
  axiom inventory.
\end{proof}

\chapter{Notes on the formalisation}
\label{chap:notes}

\section{Corrections to the paper}
\label{sec:corrections}

Formalising surfaced five places where the written text needs correcting.
None of them affects the results.

\aux{Remark}
\begin{remark}[Lemma~19: the $\lfloor m \rfloor + 1$ distinct actors]
  \label{note-lem19}
  The proof writes ``Therefore by (25), we obtain a sequence of
  $\lfloor m \rfloor + 1$ distinct actors'' without giving the construction.  One
  that works: read the actors $A_{\tau(u)-1}, A_{\tau(u)-2}, \dots$ backwards in
  time.  By (25) consecutive rows differ by exactly one expression, so their
  pressures for $O_{\tau(u)}$ form a walk that starts at $0$, ends at $m$, and
  moves by $+1$ on an expression of $O_{\tau(u)}$ and by $-\pressone$ otherwise.
  For each $j$ take the \emph{first} time the walk reaches level $\ge j + r$: a
  first passage can only happen on an up-step, of size $1$, so the walk lands
  strictly below $j+1+r$.  The $\lfloor m \rfloor + 1$ first-passage times are
  therefore distinct and name $\lfloor m \rfloor + 1$ distinct actors.

  Separately, \ref{def5-favouring} requires $n(u) = \lfloor m \rfloor \ge 1$,
  that is $m \ge 1$.  The proof rules out $m = 0$ only through the case
  $\tau(u) = 2$, but $m = 0$ forces $\Uskel{u}_{\tau(u)-1} = 0$ for any
  $\tau(u)$, so the degenerate branch has to be taken on $m = 0$ rather than on
  $\tau(u) = 2$.  The argument given for $\tau(u) = 2$ covers it unchanged.
\end{remark}

\aux{Remark}
\begin{remark}[Lemma~20: the induction invariant is not preserved]
  \label{note-lem20}
  The induction carries, for every $k$,
  $\Uskel{}_k(a,p) \le n(u) + r - \tfrac{k+1}{M-1}$ for all $a$ and all
  $p \ne o$.  The actor that expresses at step $k$ has its whole row reset to
  $0$, so its entries for $p \ne o$ are exactly $0$.  At the terminal
  $k = (M-1)(n(u)+1) - 1$ the right-hand side is $r - 1 \in [-1, 0)$, so the
  invariant asserts $0 \le r - 1 < 0$ for that actor.

  The conclusion survives: membership in $C^o$ only needs the off-columns to be
  $\le 0$, and $0$ qualifies.  Replacing the bound by
  $\max(0,\, n(u) + r - \tfrac{k+1}{M-1})$ gives an invariant that \emph{is}
  preserved — a reset row satisfies it outright, a listening row decreases by
  $\pressone$ — and still yields $\le 0$ at the terminal $k$.

  A second point in the same induction: it also asserts that $n(u)$ witness
  actors exist at every $k$.  When the maximising actor is itself a witness —
  the typical case — expressing resets it and only $n(u)-1$ survive.  The
  missing witness is replenished by the actors that have just expressed, which
  accumulate pressure for $o$ and form a staircase; that is the mechanism of
  \ref{aux-consensus-ladder}, but it is not the argument written in the proof.
\end{remark}

\aux{Remark}
\begin{remark}[Definition~4 needs a sign condition]
  \label{note-def4}
  \uses{def4}
  Definition~4 reads $0 = u(a_1,o) < u(a_2,o) < \dots < u(a_N,o)$, which says
  two things: the values are pairwise distinct, \emph{and} the smallest is $0$.
  A first formalisation of $\steep^o$ recorded only the first, together with
  ``some value is $0$'', which is strictly weaker; and the missing sign condition
  is not recoverable from the others, since the vanishing row sums hold
  identically once the relation between the columns does, whatever the sign of
  $u(a,o)$.

  \texttt{SocialNetwork.IsSteepLadder} therefore carries a \texttt{nonneg}
  field.  It costs nothing elsewhere: the only place that builds a steep ladder
  from scratch is \ref{rem5}, and the non-negativity of a ladder supplies it.
  Without it \ref{rem5-stable} is false.
\end{remark}

\aux{Remark}
\begin{remark}[Proposition~22 does not follow from Proposition~6]
  \label{note-prop22}
  \uses{rem7, prop6}
  Appendix~C says only that ``the proof of Propositions 22, 23 and 24 follows
  as the proofs of Propositions 6, 7 and 8''.  For Proposition~24 that is
  right.  For Proposition~22 the transported argument does not reach the
  stated bound, and the gap is exactly the slack $1/(2\gamma)$ that
  distinguishes $\tilde{\xi}$ from $\xi$.

  Proposition~6 runs in two cases.  If the first $N$ expressions come from $N$
  distinct actors then every row was reset, so every entry is at most $N-1$;
  that case transports unchanged, with $n_a$ in place of
  $\|U(a,\cdot)\|_\infty$.  If some actor expresses at steps $j < k < N$, then
  at step $k$ it has heard $k - j - 1$ expressions, and greediness makes its
  entry the maximum of the whole matrix, so every entry is at most $k-j-1$;
  over the remaining $N-k$ steps each entry grows by at most one, giving
  $(k-j-1) + (N-k) = N-j-1 \le N-1$.

  Under $\tilde{\xi}$ the expressed pair is only within $1/(2\gamma)$ of the
  maximum, so the second case gives
  \[
    U^{\alpha,u}_{T_N}(a,p) \;<\; N - j - 1 + \tfrac{1}{2\gamma}
    \;\le\; N - 1 + \tfrac{1}{2\gamma},
  \]
  which is $< N$ only when $\gamma \ge 1/2$.  In the regime of Appendix~C,
  $\gamma = 1/(M-1) - \alpha$ with $\alpha > 0$, so $\gamma < 1/(M-1)$: for
  every $M \ge 3$ the transported bound never gives $N$.

  The statement is not obviously false --- the slack is self-correcting, since
  $u(a,p) \le n_a$ means a large maximum forces an actor that has heard a lot
  to express, and expressing resets it.  A run started from
  $n = (0,1,2)$ with $M = N = 3$, $\gamma = 1/4$ does climb to $3$ after two
  expressions and is then pushed back to $2$ by the third, because the only
  pair within $1/(2\gamma) = 2$ of the maximum belongs to the actor carrying
  it.  What is missing is a proof that uses that feedback; the chain of
  Proposition~6, which never looks at the steps after $k$, cannot supply one.
  Either such a proof, or a weaker constant such as $N + 1/(2\gamma)$ carried
  through Propositions~23 and~17 --- both only need a bound depending on
  $\alpha, M, N$ --- would settle it.
\end{remark}

\aux{Remark}
\begin{remark}[Equation~(6) is not stable]
  \label{note-eq6}
  \uses{eq6}
  Equation~(6) pins down $\Salpha$ by two conditions on the memory profile: some
  actor has heard nothing, and $\sum_a n_a \ge N(N-1)/2$.  Remark~1 asserts that
  $\Salpha$ is stable under every $\expressa{a}{o}$.

  \textbf{The second condition is not stable.}  Expressing sets the expresser's
  $n_a$ to $0$ and adds one to every other, so the sum changes by
  $(N-1) - n_a$, which is negative once the expressing actor has heard more than
  $N-1$ expressions.  For $N = 3$, where the threshold is $3$, the profile
  $n = (0,0,3)$ has sum $3$ and satisfies the condition; expressing from the
  third actor gives $n = (1,1,0)$, of sum $2$.

  The justification the paper gives proves something stronger which \emph{is}
  stable: the actors cannot express simultaneously, so the $n_a$ are the times
  elapsed since $N$ distinct moments and are therefore \textbf{pairwise
  distinct}.  Expressing replaces $\{n_b : b \ne a\}$ by
  $\{n_b + 1 : b \ne a\}$, all at least one, and adds $n_a \mapsto 0$, so
  distinctness survives.  Distinctness together with a null row gives
  $\sum_a n_a \ge 0 + 1 + \dots + (N-1)$, so nothing is lost.

  \texttt{SocialNetwork.Bias.IsBiasedState} therefore takes distinctness as its
  second field, and the paper's inequality is derived from it.  Every state the
  process actually visits has distinct $n_a$, which is the regime Remark~1 is
  describing.
\end{remark}

\section{What is missing, and why}
\label{sec:missing}

The unproved statements fall into four groups, and it is worth keeping them
apart because they are unproved for four different reasons: the library Lean
is built on has no theory of the object, the paper's own proof does not
compose, the statement is a citation rather than a result of this paper, or it
is simply not done yet.  \texttt{STATUS.md} at the root of the repository is
the generated index of all four, statement by statement; the paragraphs below
are the argument behind it.

\subsection*{Blocked on Mathlib}

\ref{thm1-1}, \ref{thm1-2} and its skeleton half \ref{thm1-2-skeleton},
\ref{thm2}, \ref{thm3}, \ref{thm4} and their auxiliaries.  Verified against the pinned revision
(Mathlib \texttt{v4.33.0} = \texttt{db584cd6}), the following do not exist:

\begin{itemize}
  \item \textbf{Doeblin's condition $\Rightarrow$ a unique invariant measure}:
    nothing.  \texttt{Kernel.Invariant} is defined, together with three lemmas,
    and \emph{no other file in Mathlib uses it}.  \texttt{Kernel.IsIrreducible}
    is the Meyn--Tweedie definition, two trivial instances and one monotonicity
    lemma.
  \item \textbf{Recurrence for chains, and return times}: nothing.  Poincaré
    recurrence exists for a measure-preserving \emph{map}, which does not
    transport to a kernel. (Kac's lemma is likewise absent --- every
    \texttt{Kac} in Mathlib is a Kac--Moody algebra --- but it turned out not
    to be needed: the inequality $\muskel(u)\,\mathbb{E}[\Rskel{u}(u)] \le 1$,
    which is all \ref{prop9} uses, holds for every invariant probability
    measure and is proved outright at \texttt{SocialNetwork.kac\_tsum\_le}.)
  \item \textbf{Total-variation distance between measures}: nothing usable.
  \item \textbf{Poisson point processes}: nothing.  Mathlib has the Poisson
    \emph{distribution} on $\mathbb{N}$ and the Poisson limit theorem.
  \item \textbf{Quantitative convergence to $\mathrm{Exp}(1)$}:
    \texttt{TendstoInDistribution} makes the qualitative half expressible, but
    the Kolmogorov-type bound is absent.
\end{itemize}

\ref{thm2} and~\ref{thm3} live, in the paper's own architecture, on the
skeleton; they are blocked by the discrete-time items, not by the
continuous-time ones.  Only \ref{thm1-1} and~\ref{thm16} need Poisson point
processes.  The shortest path is therefore Doeblin alone: \ref{prop9} is
already proved from \ref{prop7} and~\ref{prop8}, \ref{rem5-eta} and Kac's
inequality, and \ref{thm2} follows from it by \ref{eq13} once $\muskel$ exists.

\subsection*{Blocked on the paper}

\ref{lem19}, \ref{lem20} and~\ref{aux-greedy-avoids}, and \ref{prop22}; through
them \ref{prop7}, \ref{prop9} and~\ref{prop23}.  See \ref{note-lem19},
\ref{note-lem20} and~\ref{note-prop22}, and the node of
\ref{aux-greedy-avoids}: in each case the written proof does not compose, and
the repair changes the argument rather than its presentation, so it is recorded
rather than guessed at in Lean. \ref{prop7} is nonetheless \emph{proved}, along
a route that does not need the step the written proof asserts --- the one
departure of that kind in this repository, and the reason is written at its
node.

\subsection*{Cited from outside the paper}

Four statements are not results of arXiv:2607.19651 at all, and nothing in this
library can discharge them.

\ref{prop12} and its twin for the biased process, \ref{prop12-biased}, are
Theorem~5.3 of [LM22].  They
are declared as \texttt{axiom}s rather than left carrying a \texttt{sorry},
because a \texttt{sorry} would put them in the inventory of work someone could
pick up, and they are not that.  Two are needed rather than one: stated
abstractly, over an arbitrary family of measures and an arbitrary hitting time,
the statement is \emph{inconsistent} --- the zero measure with an empty ladder
set satisfies the hypotheses vacuously and falsifies the conclusion at $t = 0$
--- so it has to be attached to a concrete process, and there are two
processes.  They are the same citation applied twice, not two assumptions.

\ref{aux-hitting-rate} and~\ref{eq19} are the two inequalities the paper
\emph{displays inside proofs} and never states. \ref{lem13} uses them, and it
cannot use the numbered statements they are attributed to --- Theorem~2.2 and
\ref{cor11} --- because those are limits, and a limit has thrown the rate away.
They are transcribed verbatim and carry a \texttt{sorry}; whether either should
become a numbered statement of the paper is a question for the authors.

\subsection*{Routine, and now done}

Seven statements were listed here as neither deep nor blocked.  They are now
proved, and the paragraph is kept as the record of what they cost.  The first
four fall in the third case of \ref{note-green}: the paper proves none of them,
so none of these is a check of anything it wrote.  The last three, added as the
earlier ones were cleared, do not: \ref{rem8} is argued in the paper and the
Lean proof follows it, and \ref{rem5-eta} and~\ref{rem5-eta-iter} supply
arguments the paper asserts.

\begin{itemize}
  \item \texttt{SocialNetwork.Bias.sum\_ge\_of\_injective}, that $N$ distinct
    naturals sum to at least $0 + 1 + \dots + (N-1)$.  This was a
    \textbf{Mathlib gap in its own right}, and still is: Mathlib has no lemma to
    this effect and none from which it follows in one step.  The route it
    suggests --- enumerate with \texttt{Finset.orderEmbOfFin}, then use ``a
    strictly monotone $\mathrm{Fin}\ N \to \mathbb{N}$ dominates the identity''
    --- needs that last step redone by hand, \texttt{StrictMono.le\_apply} being
    stated only for endomorphisms.  Induction on the largest element
    (\texttt{Finset.induction\_on\_max}) sidesteps it: adjoining a new maximum
    $a$ to $s$ adds $a$ to the sum and $\#s$ to the bound, and $a \ge \#s$
    because $s$ is contained in \texttt{range a}.  The general form is
    \texttt{SocialNetwork.Bias.sum\_range\_card\_le\_sum}.
  \item The two measurability lemmas for the biased greedy events,
    \texttt{measurableSet\_nearGreedyEvents} and
    \texttt{measurableSet\_biasedGreedyEvents}.  Both are the cylinder argument
    already used for the unbiased model, transported through
    \texttt{SocialNetwork.Bias.ofHistoryPath}.
  \item \texttt{SocialNetwork.measurable\_process}.  The characterisation is
    \texttt{jumpCount\_eq\_iff}: for $k \ne 0$, \texttt{jumpCount w t} equals $k$
    exactly when $T_k \le t$ and no later jump time is $\le t$.  The level set at
    $0$ is then the complement of the others, and \texttt{jumpCount} is
    measurable into $\mathbb{N}$; gluing it to
    \texttt{measurable\_state\_ofStepPath} along that countable index gives the
    result.
  \item \ref{rem8}.  The pigeonhole step was already there; what
    was left was the identity
    $c_p(1+\gamma) - \gamma n_a = c_p (M-1)\alpha + \gamma (M c_p - n_a)$,
    \texttt{SocialNetwork.Bias.pressure\_eq\_of\_bias}, after which both terms
    are visibly non-negative.  Cheap, as advertised.
  \item \ref{rem5-eta}, the bound $\eta$ of Remark~5.  Not cheap, and the
    one item on this list that checks something the paper wrote down without
    argument.  Two thirds of it is bookkeeping --- the positive entries of a
    steep ladder are exactly the non-minimal entries of the supported column,
    the rest are non-positive and there are at most $MN$ of them, and the
    fraction $S/(S+T)$ is monotone in the right directions.  The remaining third
    is the comparison itself, and it needs the monotone form of the same Mathlib
    gap as \texttt{sum\_range\_card\_le\_sum}: $N-1$ distinct naturals $\ge 1$
    dominate $1, 2, \dots, N-1$ term by term under any increasing $f$, which is
    \texttt{SocialNetwork.sum\_range\_add\_le\_sum}, proved by the same induction
    on the largest element.
  \item \ref{rem5-eta-iter}, the iterate $\eta^m$.  As advertised: the
    induction of \ref{prop8} along the Ionescu--Tulcea kernels, transcribed
    against the one-step bound of \ref{rem5-eta}, with the event that every
    expression so far was made from a positive entry in place of the greedy
    one.  The one real difference is that $\eta$, unlike $\zeta_\beta$, is not
    uniform in the matrix, so the finite-horizon step has to carry the
    steepness of the state the history reaches.
\end{itemize}

Nothing is left on this list, and the one item that outlived it,
\ref{aux-measurable-hitting}, is now proved as well; the subsection after next
records what its diagnosis got wrong.  What is unproved here is blocked on
Mathlib, blocked on the paper, cited from outside it, or --- for \ref{rem6}
alone --- simply not done.

\subsection*{Was here, and is now closed}

\textbf{Kac's lemma} stood in the list above as the second gap in Mathlib, and
\ref{prop9}, \ref{cor10} and three results below them were recorded as waiting
on it.  They were not.  What the paper's proof uses is the inequality
$\muskel(u)\,\mathbb{E}[\Rskel{u}(u)] \le 1$, not the identity, and the
inequality holds for \emph{every} invariant probability measure --- the
identity is what needs irreducibility, and the counterexample is two absorbing
states with $\muskel = (\tfrac12, \tfrac12)$, where the return time is $1$ and
not $2$.  The inequality is proved at \texttt{SocialNetwork.kac\_tsum\_le}, for
a Markov kernel on a countable space, from \texttt{Kernel.Invariant} and
nothing else, by decomposing the event of visiting $u$ before time $m$ over the
last such visit.  The lesson is the same as the one below: a citation is not an
obstruction until one has checked which half of it the proof needs.

\ref{aux-measurable-hitting} stood in this chapter as a fourth kind of
obstruction: not a gap in Mathlib, not a gap in the paper, not a citation, but
a statement whose Lean formulation had still to be settled.  It is now
\emph{proved}, and the diagnosis that put it here was wrong.

What was written was that reducing the infimum over $\{t \ge 0\}$ to a
countable one needs right-continuity of $t \mapsto \Uproc{u}_t(\omega)$, that
right-continuity holds only almost surely, and that the statement therefore
wanted an almost-sure formulation or an argument through the null set.  The
first clause is the error: right-continuity is \emph{a} route to the reduction,
not the only one, and the reduction holds pointwise.  The argument that works
is at \ref{aux-measurable-hitting}; it reads the level sets of the jump count
directly, and never mentions the holding times, so it costs nothing to
transport to the biased process.

The paragraph is kept because the mistake is instructive.  A property that
fails on a null set is a reason to look for a different argument before it is a
reason to weaken the statement, and an obstruction recorded in this chapter is
a claim like any other --- one that the repository owes the reader a proof or
a retraction of.

\section{What the formalised proofs check}
\label{sec:audit}

Every proof currently formalised, audited against the paper on the question of
\ref{note-green}: does it follow what the paper wrote?  Thirty-eight proofs
carry \texttt{\textbackslash leanok}; the audit below also covers the bound
$\tau(u) \le N+1$, which lives inside a definition rather than at a node of its
own, so it counts thirty-nine entries.

\medskip \noindent\textbf{Follow the paper's proof} (nineteen). \ref{prop5},
\ref{prop6}, \ref{prop7}, \ref{prop9}, \ref{prop21}, \ref{prop17}, \ref{prop24}
and~\ref{prop18}, part~1 of \ref{thm4}, the gap estimate of \ref{aux-gap}, the
bound of \ref{prop8}, \ref{lem13}, \ref{cor15} and~\ref{cor30},
\ref{thm3} and~\ref{thm31}, and \ref{rem1}, \ref{rem4} and~\ref{rem8}.  These
are the statements in the paper whose proofs are both given there and
formalised here, and each Lean proof reproduces the written argument.  Four
places carry a note: the passage from the repeat time to the global bound in
\ref{prop6}, which the paper compresses into ``this implies'', is carried out
explicitly; the iteration in \ref{prop8} goes through the Ionescu--Tulcea
kernels rather than the conditioning of~(10), Mathlib offering nothing of that
shape; \ref{rem4} names Bernoulli's inequality as the route and the degenerate
case $MNe^{-\beta/(M-1)} \ge 1$ is supplied; and \ref{rem8} keeps the count
$c_p$ the pigeonhole returns instead of rounding it down to the paper's $k =
\lceil n_a/M \rceil$, which only makes the intermediate bound larger, and
supplies $k \ge 1$.

Three more carry notes. \ref{prop21} is the proof of Proposition~5 run on $n_a$
rather than on $\|U(a,\cdot)\|_\infty$, the passage between them being $u(a,p)
\le n_a$, which equation~(6) gives but does not state. \ref{thm3} is
Section~5.3's verification of \eqref{eq:h15}--\eqref{eq:h18}, with three steps
supplied: the inclusion $\ladder \subseteq \ladder^o \cup C^{-o}$, the supremum
of~(20), and the threshold $\beta_1$; and its \ref{prop12} is an axiom, not a
theorem. \ref{cor15} is the paper's rearrangement, but needs $\ladder^o \ne
\emptyset$, which is \ref{aux-ladder-nonempty}; \ref{cor30} is that same
rearrangement with $\pressone$ replaced by $\tfrac{1}{2\gamma}$, and needs the
same witness, \ref{aux-biased-ladder-nonempty}. \ref{lem13} is the paper's case
split, resting on two inequalities the paper displays inside proofs rather than
states --- \ref{aux-hitting-rate} and~\ref{eq19} --- and supplying the
arithmetic of ``putting the inequalities above together''. \ref{thm31} is the
same verification as \ref{thm3} with $\pressone$ replaced by
$\tfrac{1}{2\gamma}$, which is all Appendix~C asks for; it names none of the
constants, so the four they take here are supplied, and it rests on a
\emph{second} axiom, the biased twin \ref{prop12-biased}, since one cannot serve
both models.
\ref{prop7} is the paper's three-stage assembly with its arithmetic unchanged,
but it applies \ref{lem20} at $\tau(u)$ rather than at $N+1$: the written route
needs $\bigcup_o S^o$ to be stable under a greedy expression, which the paper
does not establish, and starting where \ref{lem19} actually lands removes the
need.  It is the one place where a formalised proof takes a different route
through the paper's own lemmas, and the reason is written out at the node.
\ref{prop17} and~\ref{prop24} are the arguments of \ref{prop6} and~\ref{prop8}
that Appendix~C invokes; the iteration common to the two is run once in Lean,
for an arbitrary choice of admissible pairs at each profile, since that is all
either uses. \ref{prop18} is the paper's decomposition~(21), its bound~(22),
its passage through $\ln(1+x) \ge x/(1+x)$ and its event $E_\varepsilon^k$, in
that order, with the strong law of large numbers taken from Mathlib as the
paper takes it from the literature.  Two steps are supplied: the paper's
``without loss of generality $u(1,o) = 0$ for all $o$'', which is the choice of
an actor whose row is null and which equation~(6) always provides; and the
terms of index below $k$, which $E_\varepsilon^k$ does not constrain and which
the paper absorbs silently into its generic constant $C$. \ref{prop9} is the
paper's own proof --- Kac's lemma, then \ref{prop7} and \ref{rem5-eta-iter} to
bound the return time below --- with one departure and two supplied steps.  The
departure is Kac: the paper states it as the identity $1/\muskel(u) =
\mathbb{E}[\Rskel{u}(u)]$, which needs irreducibility and hence
\ref{thm1-2-skeleton}, and the proof uses only the inequality
$\muskel(u)\,\mathbb{E}[\Rskel{u}(u)] \le 1$, which holds for every invariant
probability measure and is proved here outright, so the proposition waits on
neither Doeblin's criterion nor Mathlib.  The supplied steps are the
composition of the two estimates, done on the realisation rather than by
restarting the chain, and the arithmetic of the printed constant, which needs
the regime $MNe^{-\beta/(M-1)} > 1$ treated separately.  It rests on
\ref{aux-greedy-avoids}, the ``without visiting $u$'' its proof reads off
\ref{prop7}, and adds the hypothesis $u \in \Sspace$.  Part~1 of \ref{thm4} is
the estimate the paper applies at the failure times $\eta_n$, applied at the
same time and with the same constant, but run as the single fixed-point step $q
\le (1-c) q$ rather than as a recursion over $n$; the two supplied steps ---
the Markov property itself, which the paper calls strong but applies at what is
for the skeleton a deterministic time, and the measurability of the failure
events --- are written out at the node.  Only part~1 is formalised; part~2 is
\ref{thm25}, \ref{thm27} and~\ref{thm31}, and Theorem~4 rests on them.

\medskip \noindent\textbf{Supply a step the paper asserts} (ten).
\ref{aux-state-stable}, \ref{aux-consensus-pos}, \ref{rem5-stable},
\ref{lem20-opening}, \ref{lem20-closing}, \ref{rem5-eta}
and~\ref{rem5-eta-iter}, \ref{rem5} and~\ref{rem2}, and the bound $\tau(u) \le
N+1$ of \ref{def5-repeat}.  Each corresponds to a sentence in the paper of the
form ``note that'', ``indeed'', or an unremarked clause inside a proof ---
``Note that $\ladder \subset \steep$'', ``Note that $\tau(u) \in \{2, \dots,
N+1\}$''.  All ten are true and none is deep; what they are not is checks of
anything the paper argued. \ref{rem5-eta} and~\ref{rem5-eta-iter} are the
substantial ones: the paper writes the two inequalities defining $\eta$ in a
single line and argues neither, and uses the $m$-step bound in the sketch of
\ref{prop9} without deriving it.

\medskip \noindent\textbf{No counterpart in the paper} (ten). \ref{aux-trust},
\ref{aux-ladder-consensus}, \ref{aux-state-along}, \ref{aux-consensus-ladder},
\ref{aux-ladder-nonempty}, \ref{aux-biased-ladder-nonempty},
\ref{aux-biased-ladder-consensus}, \ref{aux-shift}
and~\ref{aux-measurable-hitting} and \ref{aux-biased-ladder-steep}, together
with the rest of the plumbing of the sample space, which is not made of
numbered statements and so appears here only inside definitions and remarks.

\medskip \noindent\textbf{No proof departs from the paper for convenience, and
no proof of a statement of the paper has been invented here.} That is the
result of the audit, and the rule of \ref{note-green} is written to keep it
that way.  Where the paper's route does not close --- Lemmas~19 and~20, and
Proposition~22 --- the statement is left unproved and the obstruction is
recorded in Section~\ref{sec:corrections} above, rather than repaired by an
argument the authors have not seen.  Where the route closes but leaves the
numbered statements --- Lemma~13, whose proof reaches into the proof of
Theorem~2 and reads Corollary~11 quantitatively --- the two displays it uses
are transcribed as statements of their own, \ref{aux-hitting-rate}
and~\ref{eq19}, and carry their own \texttt{sorry}.  And where the route needs
a step the paper does not have but a shorter route through the same lemmas does
not --- \ref{prop7} --- the shorter route is taken and the difference is
written down, rather than the missing step being assumed.
\texttt{FOR-THE-AUTHORS.md} at the root of the repository collects those, and
the decisions they call for, in one place.

\aux{Remark}
\begin{remark}[``By definition'' at the end of Proposition~7]
  \label{note-prop7}
  One entry above deserves singling out.  The proof of Proposition~7 ends:

  \begin{quote}
    To conclude, we consider that by definition, if $v \in C^o$ and event
    $\bigcap_{j=1}^N \xi_j^u$ occurs, then $\Uskel{u}_N \in \ladder^o$.
  \end{quote}

  It is not by definition.  Getting from a consensus state to the exact
  staircase of \ref{def1} takes three stages, none of them formal: every one of
  the $N$ greedy expressions expresses $o$, and this has to be propagated,
  since it is stability of $C^o$ under $\express{a}{o}$ that keeps it true; no
  actor expresses twice, which is the argument of \ref{prop5} run again against
  the growth of the $o$-column; and the actor that expressed at step $i$ ends
  at exactly $N-1-i$ on column $o$, so that $i \mapsto N-1-i$ delivers
  \ref{def1} as a bijection.  In Lean this is a file of its own,
  \texttt{SocialNetwork/Consensus.lean}.

  Nothing is wrong with the claim --- unlike \ref{lem19} and~\ref{lem20}, it is
  true as written and now proved.  What is misleading is ``by definition'',
  which invites a reader to check nothing.  This is the mechanism behind the
  paper's fast consensus formation, and it earns an argument.
\end{remark}

\section{Where to go next}

In rough order of cost.

\begin{enumerate}
  \item \ref{rem6}, the one numbered statement of the paper with no Lean
    counterpart.  It is what is left that is nobody's fault, now that
    \ref{aux-measurable-hitting}, \ref{prop18}, \ref{thm4} part~1, \ref{prop9}
    and \ref{cor10} are closed --- and even it is only half free: stating it
    costs nothing, but proving it goes through \ref{cor11}, which waits on
    Doeblin below.
  \item \ref{lem19} and~\ref{lem20}, with the repairs of
    \ref{note-lem19} and~\ref{note-lem20}, and \ref{aux-greedy-avoids}, the
    step the proof of \ref{prop9} asserts.  Worth settling with the authors
    first, since the repairs change the written proofs. \ref{prop7} then
    follows by assembling the first two with \ref{aux-consensus-ladder}, and
    \ref{prop9} is proved outright once the third is.
  \item \textbf{Doeblin's condition $\Rightarrow$ a unique invariant measure for
    a countable-state Markov kernel.} The keystone: \ref{thm1-2} and its
    skeleton half \ref{thm1-2-skeleton}, and through them \ref{thm2}, \ref{thm3},
    \ref{thm25}, \ref{thm27} and \ref{thm31}, all wait on it.  It is
    self-contained, of independent value, and depends on nothing else in this
    list.
  \item \textbf{Poisson point processes} and the superposition comparison, for
    \ref{thm1-1} and~\ref{thm16}.  The largest item, and the one the fewest
    results depend on.
\end{enumerate}

Items 3 and 4 are contributions to Mathlib rather than to this repository.
The engineering companion to this blueprint, \texttt{blueprint/blueprint.md},
records exactly what the library does and does not provide, with names, files
and line numbers, checked against the pinned revision.
